% Final report submitted to the LMS as G3-report.pdf (assignment section 6.3).
% The three part chapters reuse exactly the same bodies as the per-part reports.
\documentclass[11pt,a4paper,openany]{report}
\usepackage{g3report}
% Generated by tools/build_reports.py — do not edit.
\newcommand{\grpDate}{October 2026}
\newcommand{\grpLecturer}{Dr. Lê Thành Sách}
\newcommand{\grpPages}{\url{https://trungly-hcmut.github.io/CO5177_DataVisualization_Assignment/}}
\newcommand{\grpRepo}{\url{https://github.com/trungly-hcmut/CO5177_DataVisualization_Assignment}}
\newcommand{\grpMemberRows}{1 & Ly Minh Trung & 2570349 & Tabular and time-series analysis; website design and deployment \\ 2 & Dinh Truong Tue Linh & 2570441 & Text analysis (Vietnamese news classification) \\ 3 & Nguyen Hoang Nam & 2570261 & Text analysis (Vietnamese news classification) \\}

% Generated by tools/build_reports.py — do not edit.
\newcommand{\TBPage}{\url{https://trungly-hcmut.github.io/CO5177_DataVisualization_Assignment/tabular/}}
\newcommand{\TBNotebook}{\url{https://github.com/trungly-hcmut/CO5177_DataVisualization_Assignment/blob/main/tabular/G3_Tabular_HotelBooking.ipynb}}
\newcommand{\TBVideo}{\textcolor{muted}{to be published}}
\newcommand{\TBLinks}{\href{https://trungly-hcmut.github.io/CO5177_DataVisualization_Assignment/tabular/}{Project page}\quad·\quad\href{https://github.com/trungly-hcmut/CO5177_DataVisualization_Assignment/blob/main/tabular/G3_Tabular_HotelBooking.ipynb}{Notebook on GitHub}\quad·\quad Video: \textcolor{muted}{to be published}}
\newcommand{\TBRowsRaw}{119,390}
\newcommand{\TBColsRaw}{32}
\newcommand{\TBRowsClean}{85,982}
\newcommand{\TBKeptPct}{72.0\%}
\newcommand{\TBCancelRaw}{37.0\%}
\newcommand{\TBCancelClean}{27.6\%}
\newcommand{\TBNotCancelRaw}{63.0\%}
\newcommand{\TBNotCancelClean}{72.4\%}
\newcommand{\TBDupRows}{33,241}
\newcommand{\TBDupPct}{27.8\%}
\newcommand{\TBDupRemoved}{33,226}
\newcommand{\TBZeroGuests}{180}
\newcommand{\TBAdrBad}{2}
\newcommand{\TBZeroNights}{715}
\newcommand{\TBNFeatures}{80}
\newcommand{\TBNNumeric}{26}
\newcommand{\TBNOneHot}{54}
\newcommand{\TBTrain}{60,187}
\newcommand{\TBVal}{12,897}
\newcommand{\TBTest}{12,898}
\newcommand{\TBNrBefore}{14,587}
\newcommand{\TBNrAfter}{1,010}
\newcommand{\TBNrDupPct}{93\%}
\newcommand{\TBLeadShort}{9.6\%}
\newcommand{\TBLeadLong}{67.7\%}
\newcommand{\TBDepNR}{99.4\%}
\newcommand{\TBDepNo}{28.4\%}
\newcommand{\TBSpecialZero}{47.7\%}
\newcommand{\TBSpecialThree}{16.8\%}
\newcommand{\TBSegGroups}{61\%}
\newcommand{\TBSegDirect}{15\%}
\newcommand{\TBSegCorporate}{19\%}
\newcommand{\TBMissCompany}{94.31\%}
\newcommand{\TBMissAgent}{13.69\%}
\newcommand{\TBMissCountry}{0.41\%}
\newcommand{\TBMissMeal}{0.98\%}
\newcommand{\TBSkewLeadBefore}{1.43}
\newcommand{\TBSkewLeadAfter}{−0.71}
\newcommand{\TBSkewAdrBefore}{0.88}
\newcommand{\TBSkewAdrAfter}{0.76}
\newcommand{\TBRhoLead}{+0.32}
\newcommand{\TBRhoPrevCancel}{+0.27}
\newcommand{\TBRhoSpecial}{−0.26}
\newcommand{\TBRhoParking}{−0.20}
\newcommand{\TBValDummyAcc}{0.724}
\newcommand{\TBValDummyF}{0.000}
\newcommand{\TBValDummyPR}{0.276}
\newcommand{\TBValDummyROC}{0.500}
\newcommand{\TBValLrAcc}{0.795}
\newcommand{\TBValLrF}{0.570}
\newcommand{\TBValLrPR}{0.684}
\newcommand{\TBValLrROC}{0.847}
\newcommand{\TBValDtAcc}{0.807}
\newcommand{\TBValDtF}{0.629}
\newcommand{\TBValDtPR}{0.694}
\newcommand{\TBValDtROC}{0.858}
\newcommand{\TBValRfAcc}{0.827}
\newcommand{\TBValRfF}{0.646}
\newcommand{\TBValRfPR}{0.756}
\newcommand{\TBValRfROC}{0.888}
\newcommand{\TBValHgbAcc}{0.826}
\newcommand{\TBValHgbF}{0.653}
\newcommand{\TBValHgbPR}{0.753}
\newcommand{\TBValHgbROC}{0.887}
\newcommand{\TBValTunedAcc}{0.828}
\newcommand{\TBValTunedF}{0.662}
\newcommand{\TBValTunedPR}{0.768}
\newcommand{\TBValTunedROC}{0.894}
\newcommand{\TBTuneCV}{0.766}
\newcommand{\TBTuneParams}{learning rate 0.05, 700 boosting iterations, 127 leaf nodes, at least 10 samples per leaf, L2 penalty 0.1}
\newcommand{\TBLeakF}{1.000}
\newcommand{\TBNodupF}{0.793}
\newcommand{\TBNodupPR}{0.904}
\newcommand{\TBCorrectF}{0.653}
\newcommand{\TBCorrectPR}{0.753}
\newcommand{\TBNodupGainF}{+0.14}
\newcommand{\TBNodupGainPR}{+0.15}
\newcommand{\TBThr}{0.33}
\newcommand{\TBImbDefP}{0.721}
\newcommand{\TBImbDefR}{0.613}
\newcommand{\TBImbDefF}{0.662}
\newcommand{\TBImbDefPR}{0.768}
\newcommand{\TBImbBalP}{0.613}
\newcommand{\TBImbBalR}{0.814}
\newcommand{\TBImbBalF}{0.700}
\newcommand{\TBImbBalPR}{0.768}
\newcommand{\TBImbThrP}{0.632}
\newcommand{\TBImbThrR}{0.784}
\newcommand{\TBImbThrF}{0.700}
\newcommand{\TBImbThrPR}{0.768}
\newcommand{\TBTestLrAcc}{0.800}
\newcommand{\TBTestLrP}{0.689}
\newcommand{\TBTestLrR}{0.501}
\newcommand{\TBTestLrF}{0.580}
\newcommand{\TBTestLrROC}{0.853}
\newcommand{\TBTestLrPR}{0.694}
\newcommand{\TBTestFinAcc}{0.821}
\newcommand{\TBTestFinP}{0.644}
\newcommand{\TBTestFinR}{0.786}
\newcommand{\TBTestFinF}{0.708}
\newcommand{\TBTestFinROC}{0.898}
\newcommand{\TBTestFinPR}{0.776}
\newcommand{\TBGainF}{+0.13}
\newcommand{\TBGainPR}{+0.08}
\newcommand{\TBGainROC}{+0.04}
\newcommand{\TBTN}{7,798}
\newcommand{\TBFP}{1,542}
\newcommand{\TBFN}{763}
\newcommand{\TBTP}{2,795}
\newcommand{\TBFPRate}{16.5\%}
\newcommand{\TBPositives}{3,558}
\newcommand{\TBImpLead}{0.147}
\newcommand{\TBImpCountry}{0.133}
\newcommand{\TBImpSegment}{0.108}
\newcommand{\TBImpSpecial}{0.104}
\newcommand{\TBImpDeposit}{0.027}
\newcommand{\TBImpDepositRank}{9}
\newcommand{\TBSegOnlineErr}{22.6\%}
\newcommand{\TBSegOnlineCancel}{35.8\%}
\newcommand{\TBSegDirectRecall}{0.41}
\newcommand{\TBSegCorpRecall}{0.42}
\newcommand{\TBSegGroupsRecall}{0.92}
\newcommand{\TBSegOfflineRecall}{0.87}

% Generated by tools/build_reports.py — do not edit.
\newcommand{\TXPage}{\url{https://trungly-hcmut.github.io/CO5177_DataVisualization_Assignment/text/}}
\newcommand{\TXNotebook}{\url{https://github.com/trungly-hcmut/CO5177_DataVisualization_Assignment/blob/main/text/G3_Text_VnExpressNews.ipynb}}
\newcommand{\TXVideo}{\textcolor{muted}{to be published}}
\newcommand{\TXLinks}{\href{https://trungly-hcmut.github.io/CO5177_DataVisualization_Assignment/text/}{Project page}\quad·\quad\href{https://github.com/trungly-hcmut/CO5177_DataVisualization_Assignment/blob/main/text/G3_Text_VnExpressNews.ipynb}{Notebook on GitHub}\quad·\quad Video: \textcolor{muted}{to be published}}
\newcommand{\TXRowsRaw}{3,780}
\newcommand{\TXRowsClean}{3,769}
\newcommand{\TXRemoved}{11}
\newcommand{\TXClasses}{10}
\newcommand{\TXMissingDesc}{24}
\newcommand{\TXDupWithin}{7}
\newcommand{\TXMultiSection}{2}
\newcommand{\TXNFC}{0}
\newcommand{\TXMedTitle}{12}
\newcommand{\TXMedDesc}{31}
\newcommand{\TXVocabUni}{2,950}
\newcommand{\TXVocabBi}{15,517}
\newcommand{\TXVocabChar}{19,160}
\newcommand{\TXSingleChar}{27}
\newcommand{\TXTrain}{2,638}
\newcommand{\TXVal}{565}
\newcommand{\TXTest}{566}
\newcommand{\TXClassMin}{374}
\newcommand{\TXClassMax}{391}
\newcommand{\TXGBowNb}{0.839}
\newcommand{\TXGBowLr}{0.820}
\newcommand{\TXGBowSvm}{0.819}
\newcommand{\TXGBowKnn}{0.751}
\newcommand{\TXGUniNb}{0.837}
\newcommand{\TXGUniLr}{0.848}
\newcommand{\TXGUniSvm}{0.858}
\newcommand{\TXGUniKnn}{0.837}
\newcommand{\TXGBiNb}{0.856}
\newcommand{\TXGBiLr}{0.864}
\newcommand{\TXGBiSvm}{0.863}
\newcommand{\TXGBiKnn}{0.849}
\newcommand{\TXGCharNb}{0.825}
\newcommand{\TXGCharLr}{0.847}
\newcommand{\TXGCharSvm}{0.858}
\newcommand{\TXGCharKnn}{0.819}
\newcommand{\TXGainTfidfLr}{+0.028}
\newcommand{\TXGainTfidfSvm}{+0.039}
\newcommand{\TXGainTfidfKnn}{+0.086}
\newcommand{\TXGainBiMin}{+0.005}
\newcommand{\TXGainBiMax}{+0.019}
\newcommand{\TXGainBiOverBowLr}{+0.044}
\newcommand{\TXGainBiOverBowSvm}{+0.044}
\newcommand{\TXGridBest}{TF-IDF 1–2-gram + Logistic Regression}
\newcommand{\TXGridBestF}{0.864}
\newcommand{\TXRandomVal}{0.081}
\newcommand{\TXBestC}{3}
\newcommand{\TXBestNgram}{1–3}
\newcommand{\TXBestMinDf}{1}
\newcommand{\TXCvF}{0.865}
\newcommand{\TXTunedVal}{0.857}
\newcommand{\TXHeadOnly}{0.790}
\newcommand{\TXLeadOnly}{0.844}
\newcommand{\TXBoth}{0.857}
\newcommand{\TXLeadOverHead}{+0.053}
\newcommand{\TXBothOverLead}{+0.013}
\newcommand{\TXLcFirstN}{263}
\newcommand{\TXLcFirst}{0.742}
\newcommand{\TXLcMidN}{1,319}
\newcommand{\TXLcMid}{0.845}
\newcommand{\TXLcLastN}{2,638}
\newcommand{\TXLcLast}{0.857}
\newcommand{\TXLcLastGain}{+0.012}
\newcommand{\TXRandVal}{—}
\newcommand{\TXRandTest}{0.102}
\newcommand{\TXRandAcc}{—}
\newcommand{\TXTfidfVal}{0.857}
\newcommand{\TXTfidfTest}{0.865}
\newcommand{\TXTfidfAcc}{0.866}
\newcommand{\TXEmbVal}{0.882}
\newcommand{\TXEmbTest}{0.873}
\newcommand{\TXEmbAcc}{0.873}
\newcommand{\TXHybVal}{0.888}
\newcommand{\TXHybTest}{0.892}
\newcommand{\TXHybAcc}{0.892}
\newcommand{\TXGainHyb}{+0.027}
\newcommand{\TXGainEmb}{+0.007}
\newcommand{\TXFinAcc}{0.892}
\newcommand{\TXFinF}{0.892}
\newcommand{\TXBestClass}{Sports}
\newcommand{\TXWorstClass}{Business}
\newcommand{\TXBestClassF}{0.973}
\newcommand{\TXWorstClassF}{0.814}
\newcommand{\TXSimOnePair}{Business and Sci-Tech}
\newcommand{\TXSimOne}{0.58}
\newcommand{\TXSimTwoPair}{Law and News}
\newcommand{\TXSimTwo}{0.55}
\newcommand{\TXNMis}{61}

% Generated by tools/build_reports.py — do not edit.
\newcommand{\TSPage}{\url{https://trungly-hcmut.github.io/CO5177_DataVisualization_Assignment/timeseries/}}
\newcommand{\TSNotebook}{\url{https://github.com/trungly-hcmut/CO5177_DataVisualization_Assignment/blob/main/timeseries/G3_TimeSeries_HanoiWeather.ipynb}}
\newcommand{\TSVideo}{\textcolor{muted}{to be published}}
\newcommand{\TSLinks}{\href{https://trungly-hcmut.github.io/CO5177_DataVisualization_Assignment/timeseries/}{Project page}\quad·\quad\href{https://github.com/trungly-hcmut/CO5177_DataVisualization_Assignment/blob/main/timeseries/G3_TimeSeries_HanoiWeather.ipynb}{Notebook on GitHub}\quad·\quad Video: \textcolor{muted}{to be published}}
\newcommand{\TSDays}{7,939}
\newcommand{\TSVars}{14}
\newcommand{\TSMissingDates}{0}
\newcommand{\TSStart}{1 January 2005}
\newcommand{\TSEnd}{26 September 2026}
\newcommand{\TSTrend}{+0.49}
\newcommand{\TSColdest}{6.1}
\newcommand{\TSColdestDate}{24 January 2016}
\newcommand{\TSHottest}{33.4}
\newcommand{\TSHottestDate}{4 June 2017}
\newcommand{\TSExtremes}{9}
\newcommand{\TSFirstFive}{23.50}
\newcommand{\TSLastFive}{24.13}
\newcommand{\TSDeltaFive}{+0.63}
\newcommand{\TSAnnualMin}{22.3}
\newcommand{\TSAnnualMinYear}{2011}
\newcommand{\TSAnnualMax}{24.6}
\newcommand{\TSAnnualMaxYear}{2024}
\newcommand{\TSSummer}{28.7}
\newcommand{\TSWinter}{17.4}
\newcommand{\TSSeasonGap}{11}
\newcommand{\TSJanStd}{3.40}
\newcommand{\TSJulStd}{1.62}
\newcommand{\TSDtdFeb}{1.42}
\newcommand{\TSDtdMar}{1.49}
\newcommand{\TSDtdJul}{0.78}
\newcommand{\TSAcfYear}{0.71}
\newcommand{\TSAcfAnomOne}{0.81}
\newcommand{\TSAcfAnomSeven}{0.11}
\newcommand{\TSR}{−0.36}
\newcommand{\TSNRise}{404}
\newcommand{\TSDTRise}{−1.03}
\newcommand{\TSDTOther}{+0.15}
\newcommand{\TSPRise}{26.5\%}
\newcommand{\TSPOther}{9.6\%}
\newcommand{\TSPRatio}{2.7}
\newcommand{\TSImpSeasOne}{0.70}
\newcommand{\TSImpSeasThree}{1.13}
\newcommand{\TSImpSeasSeven}{1.63}
\newcommand{\TSImpLinOne}{0.71}
\newcommand{\TSImpLinThree}{1.14}
\newcommand{\TSImpLinSeven}{1.69}
\newcommand{\TSImpClimFourteen}{1.86}
\newcommand{\TSImpClimThirty}{1.87}
\newcommand{\TSRainOutliers}{1,001}
\newcommand{\TSRainSkew}{5.2}
\newcommand{\TSRainSkewLog}{0.95}
\newcommand{\TSRainMax}{227}
\newcommand{\TSRainMaxDate}{21 June 2010}
\newcommand{\TSTrain}{6,179}
\newcommand{\TSVal}{730}
\newcommand{\TSTest}{1,000}
\newcommand{\TSPhi}{0.81}
\newcommand{\TSCvTuned}{0.928}
\newcommand{\TSCvDefault}{0.944}
\newcommand{\TSValClim}{2.014}
\newcommand{\TSTestClim}{1.990}
\newcommand{\TSValPersist}{1.105}
\newcommand{\TSTestPersist}{1.081}
\newcommand{\TSValDamped}{1.092}
\newcommand{\TSTestDamped}{1.075}
\newcommand{\TSValRidge}{0.974}
\newcommand{\TSTestRidge}{0.970}
\newcommand{\TSValRf}{0.921}
\newcommand{\TSTestRf}{0.935}
\newcommand{\TSValLevel}{0.933}
\newcommand{\TSTestLevel}{0.932}
\newcommand{\TSValTuned}{0.922}
\newcommand{\TSTestTuned}{0.919}
\newcommand{\TSValFinal}{0.906}
\newcommand{\TSTestFinal}{0.916}
\newcommand{\TSValSkill}{18\%}
\newcommand{\TSTestSkill}{15.3\%}
\newcommand{\TSTestRmse}{1.226}
\newcommand{\TSTestRtwo}{0.935}
\newcommand{\TSMoFebModel}{1.16}
\newcommand{\TSMoFebPersist}{1.61}
\newcommand{\TSMoMarModel}{0.82}
\newcommand{\TSMoMarPersist}{1.24}
\newcommand{\TSMoDecModel}{0.88}
\newcommand{\TSMoDecPersist}{1.12}
\newcommand{\TSMoAprModel}{1.10}
\newcommand{\TSMoAprPersist}{1.08}
\newcommand{\TSMoNovModel}{1.12}
\newcommand{\TSMoNovPersist}{1.01}
\newcommand{\TSMonthsWorse}{2}
\newcommand{\TSResMean}{−0.27}
\newcommand{\TSResStd}{1.20}
\newcommand{\TSWorstN}{8}
\newcommand{\TSWorstCold}{7}
\newcommand{\TSWorstDropMin}{4.6}
\newcommand{\TSWorstDropMax}{8.5}
\newcommand{\TSWorstErrMin}{4.2}
\newcommand{\TSWorstErrMax}{6.7}
\newcommand{\TSUseful}{4}
\newcommand{\TSPersistCross}{3}
\newcommand{\TSPersistCrossMae}{2.161}
\newcommand{\TSClim}{1.990}
\newcommand{\TSHOne}{0.916}
\newcommand{\TSHOneBase}{1.075}
\newcommand{\TSHTwo}{1.489}
\newcommand{\TSHTwoBase}{1.603}
\newcommand{\TSHThree}{1.792}
\newcommand{\TSHThreeBase}{1.858}
\newcommand{\TSHFour}{1.946}
\newcommand{\TSHFourBase}{1.974}
\newcommand{\TSHLast}{1.989}
\newcommand{\TSImpAnom}{0.300}
\newcommand{\TSImpWind}{0.196}
\newcommand{\TSImpRain}{0.134}
\newcommand{\TSImpCloud}{0.087}
\newcommand{\TSImpDp}{0.062}

\setpartlabel{Final report}
\hypersetup{pdftitle={Group G3 — Final report — Programming Foundations for Data Analysis and Visualization}}
\setcounter{tocdepth}{1}

\begin{document}
\coverpage{Final assignment report · Group G3}
  {Three Data Types, One Workflow}
  {Hotel booking cancellations · Vietnamese news topics · Hanoi temperature forecasts}
  {\href{https://trungly-hcmut.github.io/CO5177_DataVisualization_Assignment/}{Project website}\quad·\quad
   \href{https://github.com/trungly-hcmut/CO5177_DataVisualization_Assignment}{GitHub repository}\quad·\quad
   Videos: \textcolor{muted}{to be published}}

\tableofcontents

% ------------------------------------------------------------------
\sectionsflow
\chapter{Group and project overview}\label{ov:chapter}

\section{Group information}
Group \textbf{G3} carried out the assignment of the course \emph{Nền tảng lập trình cho phân tích và trực quan dữ liệu}
(Programming Foundations for Data Analysis and Visualization), Semester 261, academic year 2026--2027, under the supervision of
\grpLecturer{} at the Faculty of Computer Science and Engineering, Ho Chi Minh City University of Technology (VNU-HCM).

\begin{table}[H]\centering\small
\caption{Members of Group G3.}\label{ov:tab-members}
\begin{tabular}{@{}clll@{}}
\toprule \hd{No.} & \hd{Full name} & \hd{Student ID} & \hd{Role} \\ \midrule
\grpMemberRows
\bottomrule
\end{tabular}
\end{table}

As a group of three, we completed the three data types the assignment requires (section 3.2): the mandatory
\textbf{tabular} and \textbf{text} parts, and \textbf{time series} as the optional third part.

\section{Deliverables}
Everything is published on the group website (GitHub Pages) at \grpPages, with the source code in \grpRepo. Each part has its own
page that shows the full executed notebook together with its report and video.

\begin{table}[H]\centering\small
\caption{Deliverables of each part (full addresses in Appendix~\ref{app:links}).}\label{ov:tab-deliverables}
\begin{tabularx}{\linewidth}{lYYYY}
\toprule \hd{Part} & \hd{Project page} & \hd{Notebook (.ipynb)} & \hd{PDF report} & \hd{Video} \\ \midrule
Tabular & \href{https://trungly-hcmut.github.io/CO5177_DataVisualization_Assignment/tabular/}{page} & \href{https://github.com/trungly-hcmut/CO5177_DataVisualization_Assignment/blob/main/tabular/G3_Tabular_HotelBooking.ipynb}{.ipynb on GitHub} & \href{https://trungly-hcmut.github.io/CO5177_DataVisualization_Assignment/tabular/report.pdf}{report.pdf} & \TBVideo \\
Text & \href{https://trungly-hcmut.github.io/CO5177_DataVisualization_Assignment/text/}{page} & \href{https://github.com/trungly-hcmut/CO5177_DataVisualization_Assignment/blob/main/text/G3_Text_VnExpressNews.ipynb}{.ipynb on GitHub} & \href{https://trungly-hcmut.github.io/CO5177_DataVisualization_Assignment/text/report.pdf}{report.pdf} & \TXVideo \\
Time series & \href{https://trungly-hcmut.github.io/CO5177_DataVisualization_Assignment/timeseries/}{page} & \href{https://github.com/trungly-hcmut/CO5177_DataVisualization_Assignment/blob/main/timeseries/G3_TimeSeries_HanoiWeather.ipynb}{.ipynb on GitHub} & \href{https://trungly-hcmut.github.io/CO5177_DataVisualization_Assignment/timeseries/report.pdf}{report.pdf} & \TSVideo \\
\bottomrule
\end{tabularx}
\end{table}

\section{The three datasets}
\begin{table}[H]\centering\small
\caption{Datasets used, one per data type.}\label{ov:tab-datasets}
\begin{tabularx}{\linewidth}{lL{3.9cm}L{3.3cm}Y}
\toprule \hd{Part} & \hd{Dataset and source} & \hd{Size} & \hd{Task} \\ \midrule
Tabular & Hotel Booking Demand (Data in Brief, 2019), two hotels in Portugal & \TBRowsRaw{} bookings × \TBColsRaw{} columns & binary classification: will the booking be cancelled? \\
Text & VnExpress headlines and leads, \textbf{crawled by the group} on 27~September~2026 & \TXRowsRaw{} articles, \TXClasses{} sections & multi-class classification: which section does the article belong to? \\
Time series & Open-Meteo Historical Weather API (ERA5 reanalysis), central Hanoi & \TSDays{} days × \TSVars{} variables (2005–2026) & regression: tomorrow's mean temperature, and 1–7 days ahead \\
\bottomrule
\end{tabularx}
\end{table}

\section{Key results}
\begin{table}[H]\centering\small
\caption{Headline result of each part on its held-out test set.}\label{ov:tab-results}
% Generated by tools/build_reports.py — do not edit.
\begin{tabular}{lL{3.0cm}L{3.2cm}lrr}
\toprule
\hd{Part} & \hd{Task} & \hd{Dataset} & \hd{Test metric} & \hd{Baseline} & \hd{Ours} \\
\midrule
Tabular & Will a hotel booking be cancelled? (binary) & Hotel Booking Demand, 119,390 bookings & F1 & 0.580 (Log. Reg.) & \textbf{0.708} \\
Text & Which of 10 sections is a news article in? & VnExpress, 3,780 articles (self-crawled) & macro-F1 & 0.865 (TF-IDF + SVM) & \textbf{0.892} \\
Time series & Hanoi's mean temperature tomorrow & Open-Meteo / ERA5, 7,939 days & MAE (°C) & 1.081 (persistence) & \textbf{0.916} \\
\bottomrule
\end{tabular}

\end{table}

\section{Shared workflow and tools}
All three parts follow the same workflow required by the assignment (section 5), and every notebook runs end-to-end on Google
Colab (\emph{Run all}) with fixed random seeds, loading its data from the repository.

\begin{figure}[H]\centering
\begin{tikzpicture}[node distance=0.36cm, box/.style={draw=bklight, fill=tint, rounded corners=2pt, font=\sffamily\footnotesize,
  text width=2.3cm, align=center, minimum height=1.35cm}, arr/.style={-{Stealth[length=2mm]}, bkdark, thick}]
\node[box] (a) {\textbf{Collect}\\download, crawl, API};
\node[box, right=of a] (b) {\textbf{Describe \& explore}\\quality checks, EDA};
\node[box, right=of b] (c) {\textbf{Prepare}\\missing, outliers, encoding, scaling, split};
\node[box, right=of c] (d) {\textbf{Model}\\baselines, scikit-learn models, tuning};
\node[box, right=of d] (e) {\textbf{Explain}\\test metrics, errors, extra experiments};
\draw[arr] (a) -- (b); \draw[arr] (b) -- (c); \draw[arr] (c) -- (d); \draw[arr] (d) -- (e);
\end{tikzpicture}
\caption{The workflow shared by the three parts.\label{ov:fig-workflow}}
\end{figure}

\textbf{Tools.} Python with NumPy and pandas (preparation), Matplotlib and Seaborn (visualisation), scikit-learn (pipelines,
models, tuning, evaluation); requests and BeautifulSoup for crawling, and the sentence-transformers library for the
pretrained text encoder. Each notebook also uses object-oriented code: custom scikit-learn transformers (\code{Winsorizer},
\code{HotelFeatureBuilder}, \code{VietnameseTextCleaner}), wrappers (\code{DeltaRegressor}) and experiment classes
(\code{ModelBenchmark}, \code{TextBenchmark}, \code{ForecastEvaluator}).

\textbf{How each finding is presented.} Every extra experiment answers one explicit question with a figure or table, numbers
to compare (before/after, model A vs B, correct vs wrong) and a conclusion in a highlighted box, as asked in section 7.1 of the
assignment. All numbers in this report are generated from the notebooks' saved results, so the text always matches the tables.

% ------------------------------------------------------------------
\sectionsonnewpage   % the part chapters mirror the standalone reports: one section per page start
\chapter{Tabular data: predicting hotel booking cancellations}\label{ch:tabular}
% Body of the Tabular report — used by tabular.tex (as an article) and G3-report.tex (as a chapter).
% Every number comes from generated/tabular.tex (built from tabular/results.json).

\section{Problem and objective}\label{tb:problem}
Hotels take bookings weeks or months in advance, yet \TBCancelRaw{} of the bookings in our dataset were cancelled.
A hotel that can tell, at booking time, which reservations are likely to be cancelled can overbook in a controlled way,
ask guests to reconfirm, or require a deposit for risky bookings.

\textbf{Objective.} Build a binary classifier that predicts \code{is\_canceled} (1 = cancelled) using only information
that is available when the booking is made, and measure honestly how well it works on bookings it has never seen.

\textbf{How success is measured.} Cancellations are the minority class, so accuracy is misleading: a model that always
answers ``not cancelled'' is already \TBNotCancelClean{} accurate on the cleaned data while catching no cancellation at all.
We therefore report \textbf{F1, recall, ROC-AUC and PR-AUC} (average precision), select models on a validation set, and
touch the test set only once.

\section{Dataset}\label{tb:dataset}
\textbf{Hotel Booking Demand} \cite{tb-antonio} contains real bookings from two hotels in Portugal — a \emph{City Hotel} in
Lisbon and a \emph{Resort Hotel} in the Algarve — with arrival dates from July 2015 to August 2017 (licence CC BY 4.0).
Each row is one booking: \TBRowsRaw{} rows and \TBColsRaw{} columns, of which 20 are numerical and 12 categorical.

\begin{table}[H]\centering\small
\caption{Column groups of the dataset.}\label{tb:tab-columns}
\begin{tabularx}{\linewidth}{L{2.5cm}YL{4.1cm}}
\toprule \hd{Group} & \hd{Columns} & \hd{Meaning} \\ \midrule
Target & \code{is\_canceled} & 1 if the booking was cancelled \\
Timing & \code{lead\_time}, \code{arrival\_date\_*} & days booked in advance; arrival date \\
Stay & \code{stays\_in\_weekend\_nights}, \code{stays\_in\_week\_nights}, \code{adults}, \code{children}, \code{babies} & nights and guests \\
Channel \& guest & \code{market\_segment}, \code{distribution\_channel}, \code{customer\_type}, \code{country}, \code{agent}, \code{company}, \code{is\_repeated\_guest} & booking source, customer type, nationality \\
History & \code{previous\_cancellations}, \code{previous\_bookings\_not\_canceled} & the guest's past behaviour \\
Service \& price & \code{meal}, \code{reserved\_room\_type}, \code{deposit\_type}, \code{adr}, \code{required\_car\_parking\_spaces}, \code{total\_of\_special\_requests}, \code{booking\_changes}, \code{days\_in\_waiting\_list} & meal plan, room, deposit, average daily rate (ADR) \\
\textbf{Leakage} & \code{reservation\_status}, \code{reservation\_status\_date}, \code{assigned\_room\_type} & known only \emph{after} arrival or cancellation — dropped \\
\bottomrule
\end{tabularx}
\end{table}

\begin{table}[H]\centering\small
\caption{Assignment requirements for tabular data (section 4.2) and how this dataset meets them.}\label{tb:tab-req}
\begin{tabularx}{\linewidth}{L{4.0cm}Yl}
\toprule \hd{Requirement} & \hd{This dataset} & \hd{Status} \\ \midrule
At least 2,000 samples & \TBRowsRaw{} bookings (\TBRowsClean{} after cleaning) & \met \\
At least 10 columns & \TBColsRaw{} columns & \met \\
Missing values & \code{company} \TBMissCompany, \code{agent} \TBMissAgent, \code{country} \TBMissCountry, \code{children}; plus hidden \code{Undefined} values & \met \\
Categorical columns (encoding) & 12, e.g. \code{hotel}, \code{meal}, \code{country}, \code{market\_segment}, \code{deposit\_type} & \met \\
Numerical columns (scaling) & 20, e.g. \code{lead\_time}, \code{adr}, \code{stays\_in\_week\_nights} & \met \\
Harder data (preferred) & imbalance (\TBNotCancelRaw{} / \TBCancelRaw{} raw, \TBNotCancelClean{} / \TBCancelClean{} after cleaning), outliers (ADR up to 5,400), \TBDupPct{} duplicates, target leakage & \met \\
\bottomrule
\end{tabularx}
\end{table}

\section{Exploratory data analysis}\label{tb:eda}
\subsection{Data quality}
Table~\ref{tb:tab-quality} lists every quality problem we found and what we did about it. Missing values are not always
errors: a missing \code{company} means ``not booked through a company'', so it becomes a binary flag rather than a deletion.
The dataset also hides missing values behind the string \code{Undefined}.

\begin{table}[H]\centering\small
\caption{Data-quality issues and the action taken for each.}\label{tb:tab-quality}
\begin{tabularx}{\linewidth}{L{3.6cm}L{3.9cm}Y}
\toprule \hd{Issue} & \hd{Affected} & \hd{Action} \\ \midrule
Missing \code{company} & \TBMissCompany{} of rows & flag \code{has\_company} (missing = no company) \\
Missing \code{agent} & \TBMissAgent{} of rows & flag \code{has\_agent} (missing = booked directly) \\
Missing \code{country}, \code{children} & \TBMissCountry{} of rows; 4 rows & most-frequent / median imputation, fitted on train \\
Hidden \code{Undefined} & \code{meal} \TBMissMeal; \code{market\_segment} 2 rows; \code{distribution\_channel} 5 rows & \code{meal}: Undefined = SC (per documentation); others → missing → most frequent \\
Duplicated rows & \TBDupRows{} (\TBDupPct) & removed before splitting \\
Invalid bookings & \TBZeroGuests{} with 0 guests; \TBAdrBad{} with ADR $<0$ or $>1{,}000$ & removed \\
0-night bookings & \TBZeroNights{} rows & kept — day use is a valid booking \\
Leakage columns & 3 columns & dropped (Section~\ref{tb:traps} shows why) \\
\bottomrule
\end{tabularx}
\end{table}

Removing duplicates changes the picture: the cancellation rate falls from \TBCancelRaw{} to \TBCancelClean{}, because most
copies were \emph{cancelled} group bookings. In particular, \TBNrDupPct{} of the \code{Non Refund} bookings were copies
(\TBNrBefore{} → \TBNrAfter). The clean data is therefore more imbalanced (\TBNotCancelClean{} / \TBCancelClean).

\subsection{Cancellation versus key features}
\fig{tab_03_target_vs_features}{Cancellation rate by lead time, deposit type, market segment and number of special requests
(raw data, groups with at least 100 bookings). The horizontal line is the overall rate.\label{tb:fig-target}}
\begin{finding}
\Q Which booking-time attributes are most strongly associated with cancellation?
\F \textbf{Lead time}: \TBLeadShort{} of bookings made at most 7 days ahead are cancelled, versus \TBLeadLong{} of bookings
made more than a year ahead. \textbf{Non Refund} deposits are cancelled \TBDepNR{} of the time (vs \TBDepNo{} for no deposit)
— counter-intuitive, and largely produced by the duplicated group blocks found above. \textbf{Groups} cancel most
(\TBSegGroups), \textbf{Direct} (\TBSegDirect) and \textbf{Corporate} (\TBSegCorporate) least. Guests with \textbf{special
requests} cancel far less (\TBSpecialZero{} with none vs \TBSpecialThree{} with three or more).
\end{finding}

\subsection{Distributions, outliers and correlation}
The numerical columns are strongly right-skewed and contain thousands of points beyond the IQR fences. Most are real
(bookings made two years ahead, expensive suites), so we do not delete them; we clip and transform them instead
(Section~\ref{tb:prep}). Spearman correlations with the target are only weak to moderate — \code{lead\_time}
\TBRhoLead, \code{previous\_cancellations} \TBRhoPrevCancel, \code{total\_of\_special\_requests} \TBRhoSpecial,
\code{required\_car\_parking\_spaces} \TBRhoParking{} — so no single column decides the outcome and a multivariate,
non-linear model is needed.

\section{Data preparation}\label{tb:prep}
All steps that learn a parameter (medians, clipping quantiles, means and standard deviations, one-hot categories) are
fitted on the training set only, inside a scikit-learn \code{Pipeline} with a \code{ColumnTransformer}. Row-level cleaning
learns nothing from the data, so it is done before splitting. The notebook implements the steps as reusable classes:
\code{HotelDataCleaner} (row cleaning with a log), \code{HotelFeatureBuilder} and \code{Winsorizer} (custom scikit-learn
transformers) and \code{ModelBenchmark} (trains and scores many models the same way).

\begin{table}[H]\centering\small
\caption{Preparation pipeline.}\label{tb:tab-pipeline}
\begin{tabularx}{\linewidth}{L{2.9cm}YL{2.2cm}}
\toprule \hd{Step} & \hd{What is done} & \hd{Fitted on} \\ \midrule
Row cleaning & drop leakage columns; remove invalid bookings; remove \TBDupRemoved{} duplicates & — (rules) \\
Feature engineering & \code{total\_nights}, \code{total\_guests}, \code{has\_children}, \code{is\_weekend\_only}, \code{has\_agent}, \code{has\_company}, \code{is\_domestic}, \code{arrival\_weekday}; cyclical \code{month\_sin}/\code{month\_cos} & — (stateless) \\
Missing values & median (numerical), most frequent (categorical) & train \\
Outliers & \code{Winsorizer}: clip to the 0.5\%–99.5\% quantiles & train \\
Skewness & \code{log1p} on count/day columns; not on \code{adr} (1.6\% zeros) & — \\
Encoding & \code{OneHotEncoder}, categories below 0.5\% merged into ``infrequent'' → \TBNOneHot{} columns & train \\
Scaling & \code{StandardScaler} on the \TBNNumeric{} numerical features & train \\
Split & stratified 70 / 15 / 15: \TBTrain{} / \TBVal{} / \TBTest{} bookings & — \\
\bottomrule
\end{tabularx}
\end{table}

\fig[0.9\linewidth]{tab_06_scaling}{Distributions of \code{lead\_time} and \code{adr} before and after preprocessing (training set).\label{tb:fig-scaling}}
\begin{finding}
\Q How does preprocessing change the two most important numerical features?
\F Both end with mean 0 and standard deviation 1. The skewness of \code{lead\_time} drops from \TBSkewLeadBefore{} to
\TBSkewLeadAfter{} and that of \code{adr} from \TBSkewAdrBefore{} to \TBSkewAdrAfter. A first version also applied
\code{log1p} to \code{adr} and made its skewness \emph{worse} (−3.6) because complimentary rooms cost 0 — which is why
\code{adr} has its own branch.
\end{finding}

\section{Modelling}\label{tb:models}
\subsection{Model comparison}
Five classifiers share the same preprocessing pipeline: a dummy baseline, Logistic Regression, a Decision Tree (depth 12),
a Random Forest (200 trees) and HistGradientBoosting (HGB).

\begin{table}[H]\centering\small
\caption{Validation scores (threshold 0.5). The tuned model is described in Section~\ref{tb:tuning}.}\label{tb:tab-val}
% Generated by tools/build_reports.py — do not edit.
\begin{tabular}{lrrrrrr}
\toprule
\hd{Model} & \hd{Accuracy} & \hd{Precision} & \hd{Recall} & \hd{F1} & \hd{ROC-AUC} & \hd{PR-AUC} \\
\midrule
Dummy (always “not cancelled”) & 0.724 & 0.000 & 0.000 & 0.000 & 0.500 & 0.276 \\
Logistic Regression & 0.795 & 0.675 & 0.493 & 0.570 & 0.847 & 0.684 \\
Decision Tree & 0.807 & 0.669 & 0.593 & 0.629 & 0.858 & 0.694 \\
Random Forest & 0.827 & 0.743 & 0.572 & 0.646 & 0.888 & 0.756 \\
HistGradientBoosting (HGB) & 0.826 & 0.723 & 0.596 & 0.653 & 0.887 & 0.753 \\
\textbf{HGB, tuned} & \textbf{0.828} & \textbf{0.721} & \textbf{0.613} & \textbf{0.662} & \textbf{0.894} & \textbf{0.768} \\
\bottomrule
\end{tabular}

\end{table}

The dummy model reaches \TBValDummyAcc{} accuracy with F1 = 0, confirming that accuracy is not informative here. The two
ensembles lead and are almost tied — HGB (F1 \TBValHgbF, PR-AUC \TBValHgbPR) and Random Forest (F1 \TBValRfF, PR-AUC
\TBValRfPR) — while Logistic Regression is last (F1 \TBValLrF, PR-AUC \TBValLrPR) despite careful scaling: the signal is
non-linear and involves interactions (a long lead time is risky only in some segments).

\subsection{Hyperparameter tuning}\label{tb:tuning}
HGB was tuned with \code{RandomizedSearchCV} (12 configurations × 3-fold stratified cross-validation on the training set,
scored by PR-AUC). The best configuration — \TBTuneParams{} — reached a cross-validated PR-AUC of \TBTuneCV{} and raised the
validation PR-AUC from \TBValHgbPR{} to \TBValTunedPR{} (F1 from \TBValHgbF{} to \TBValTunedF).

\section{Results on the test set}\label{tb:results}
The final model is the tuned HGB with the decision threshold chosen on the validation set (Section~\ref{tb:imbalance}).
It was evaluated once on the \TBTest{} test bookings.

\begin{table}[H]\centering\small
\caption{Test-set results of the baseline and the final model.}\label{tb:tab-test}
% Generated by tools/build_reports.py — do not edit.
\begin{tabular}{lrrrrrr}
\toprule
\hd{Model (test set)} & \hd{Accuracy} & \hd{Precision} & \hd{Recall} & \hd{F1} & \hd{ROC-AUC} & \hd{PR-AUC} \\
\midrule
Logistic Regression (baseline) & 0.800 & 0.689 & 0.501 & 0.580 & 0.853 & 0.694 \\
\textbf{HGB tuned, threshold 0.33} & \textbf{0.821} & \textbf{0.644} & \textbf{0.786} & \textbf{0.708} & \textbf{0.898} & \textbf{0.776} \\
\bottomrule
\end{tabular}

\end{table}

\fig{tab_10_test_eval}{Final model on the test set: confusion matrix, ROC curve and precision–recall curve, compared with
the Logistic Regression baseline.\label{tb:fig-test}}
\begin{finding}
\Q How good is the final model on bookings it has never seen?
\F F1 rises from \TBTestLrF{} (baseline) to \textbf{\TBTestFinF} (\TBGainF), ROC-AUC from \TBTestLrROC{} to
\textbf{\TBTestFinROC} and PR-AUC from \TBTestLrPR{} to \textbf{\TBTestFinPR}. Of the \TBPositives{} real cancellations
the model catches \TBTP{} (recall \TBTestFinR); in exchange \TBFP{} bookings that were kept (\TBFPRate) are flagged by mistake.
\end{finding}

\section{Extra experiments}\label{tb:extra}
\subsection{Two evaluation traps: leakage and duplicates}\label{tb:traps}
\begin{table}[H]\centering\small
\caption{Validation scores of the same HGB model under three data-handling scenarios.}\label{tb:tab-pitfalls}
% Generated by tools/build_reports.py — do not edit.
\begin{tabular}{lrrrr}
\toprule
\hd{Scenario} & \hd{Accuracy} & \hd{F1} & \hd{ROC-AUC} & \hd{PR-AUC} \\
\midrule
Leakage columns kept & 1.000 & 1.000 & 1.000 & 1.000 \\
Duplicates not removed & 0.854 & 0.793 & 0.933 & 0.904 \\
\textbf{Correct process (this report)} & \textbf{0.826} & \textbf{0.653} & \textbf{0.887} & \textbf{0.753} \\
\bottomrule
\end{tabular}

\end{table}
\fig[0.82\linewidth]{tab_08_pitfalls}{F1 and PR-AUC when the cleaning steps are skipped.\label{tb:fig-pitfalls}}
\begin{finding}
\Q If we skipped the cleaning of Section~\ref{tb:prep}, how good would the scores look?
\F Keeping the leakage columns gives \textbf{\TBLeakF{} on every metric}: the model just reads the answer from
\code{reservation\_status}. Keeping duplicates inflates F1 from \TBCorrectF{} to \TBNodupF{} (\TBNodupGainF) and PR-AUC from
\TBCorrectPR{} to \TBNodupPR{} (\TBNodupGainPR), because copies of the same booking appear in both training and validation.
The honest figures are the lower ones.
\end{finding}

\subsection{Class imbalance and the decision threshold}\label{tb:imbalance}
\begin{table}[H]\centering\small
\caption{Three ways to handle the imbalance, on the tuned HGB model (validation set).}\label{tb:tab-imbalance}
% Generated by tools/build_reports.py — do not edit.
\begin{tabular}{lrrrr}
\toprule
\hd{Strategy (validation set)} & \hd{Precision} & \hd{Recall} & \hd{F1} & \hd{PR-AUC} \\
\midrule
Default threshold 0.5 & 0.721 & 0.613 & 0.662 & 0.768 \\
\code{class\_weight=balanced} & 0.613 & 0.814 & 0.700 & 0.768 \\
\textbf{Threshold tuned for F1 (0.33)} & \textbf{0.632} & \textbf{0.784} & \textbf{0.700} & \textbf{0.768} \\
\bottomrule
\end{tabular}

\end{table}
\fig{tab_09_imbalance}{Left: precision, recall and F1 as the threshold moves. Right: the three strategies compared.\label{tb:fig-imbalance}}
\begin{finding}
\Q Does handling the imbalance help, and how?
\F At threshold 0.5 the model misses many cancellations (recall \TBImbDefR). Re-weighting the classes raises recall to
\TBImbBalR{} and lowering the threshold to \TBThr{} raises it to \TBImbThrR; F1 rises from \TBImbDefF{} to \TBImbBalF{} with
re-weighting and to \TBImbThrF{} with the tuned threshold. PR-AUC barely moves (\TBImbDefPR): handling imbalance shifts the operating point along the same
curve — trading precision for recall — rather than making the model better at ranking. We keep the tuned threshold because
it needs no retraining and can later be set from the hotel's real costs.
\end{finding}

\subsection{What the model relies on}
\fig[0.82\linewidth]{tab_11_importance}{Permutation importance on the test set: drop in PR-AUC when one original column is shuffled.\label{tb:fig-importance}}
\begin{finding}
\Q Which columns does the model depend on most?
\F \code{lead\_time} (\TBImpLead) and \code{country} (\TBImpCountry) matter most, followed by \code{market\_segment}
(\TBImpSegment) and \code{total\_of\_special\_requests} (\TBImpSpecial) — consistent with Figure~\ref{tb:fig-target}.
\code{deposit\_type} ranks only \TBImpDepositRank{}th (\TBImpDeposit): once duplicates are removed, very few Non Refund
bookings remain. Permutation importance shows what the model \emph{uses}, not what \emph{causes} a cancellation.
\end{finding}

\subsection{Where the model fails}
\begin{table}[H]\centering\small
\caption{Test errors by market segment (segments with at least 50 test bookings).}\label{tb:tab-segments}
% Generated by tools/build_reports.py — do not edit.
\begin{tabular}{lrrrr}
\toprule
\hd{Market segment} & \hd{Test bookings} & \hd{Cancel rate} & \hd{Error rate} & \hd{Recall} \\
\midrule
Online TA & 7,693 & 35.8\% & 22.6\% & 0.81 \\
Direct & 1,794 & 13.0\% & 13.9\% & 0.41 \\
Complementary & 95 & 11.6\% & 13.7\% & 0.09 \\
Corporate & 571 & 12.6\% & 10.9\% & 0.42 \\
Offline TA/TO & 2,027 & 14.4\% & 8.9\% & 0.87 \\
Groups & 686 & 27.1\% & 8.5\% & 0.92 \\
\bottomrule
\end{tabular}

\end{table}
\begin{finding}
\Q Is the error rate the same for every kind of booking?
\F No. Errors concentrate in \textbf{Online TA} (error rate \TBSegOnlineErr), the largest segment and the one with the most
cancellations (\TBSegOnlineCancel). For \textbf{Direct} and \textbf{Corporate} bookings recall is only \TBSegDirectRecall{} and
\TBSegCorpRecall: these guests rarely cancel, so the model seldom predicts a cancellation for them. Groups
(\TBSegGroupsRecall) and Offline TA/TO (\TBSegOfflineRecall) are caught well. A per-segment threshold would be a natural
next step.
\end{finding}

\section{Conclusions}\label{tb:conclusion}
\begin{enumerate}
\item \textbf{Cleaning decides whether results can be trusted.} \TBDupPct{} of the rows were duplicates and three columns
leaked the answer; ignoring them gives ``excellent'' scores (F1 \TBNodupF{}–\TBLeakF) that are simply wrong. The true
cancellation rate after cleaning is \TBCancelClean, not \TBCancelRaw.
\item \textbf{Non-linear ensembles work best.} The tuned HistGradientBoosting model beats Logistic Regression by \TBGainF{}
F1 and \TBGainPR{} PR-AUC on the test set (F1 \TBTestFinF, ROC-AUC \TBTestFinROC).
\item \textbf{The threshold matters as much as the model.} Moving it from 0.5 to \TBThr{} raises validation recall from
\TBImbDefR{} to \TBImbThrR{} at the same ranking quality.
\item \textbf{Strongest signals:} booking far ahead, nationality, sales channel, and \emph{not} making special requests.
\end{enumerate}
\textbf{Limitations.} The split is random rather than chronological, so the score on truly future bookings may be lower;
the data covers two Portuguese hotels before COVID-19; and the F1-optimal threshold assumes that a missed cancellation and a
false alarm cost the same.

\textbf{Future work.} A time-based split (train 2015–2016, test 2017); a threshold derived from real business costs;
\code{TargetEncoder} for high-cardinality \code{country} and \code{agent}; probability calibration so scores can drive
overbooking directly; per-booking explanations with SHAP.

\begin{thebibliography}{9}
\bibitem{tb-antonio} N.~Antonio, A.~de~Almeida, L.~Nunes. \emph{Hotel booking demand datasets}. Data in Brief, 22 (2019),
  41--49. \url{https://doi.org/10.1016/j.dib.2018.11.126}
\bibitem{tb-sklearn} F.~Pedregosa et al. \emph{Scikit-learn: Machine Learning in Python}. Journal of Machine Learning
  Research, 12 (2011), 2825--2830.
\bibitem{tb-nb} Group G3. \emph{Notebook} \code{G3\_Tabular\_HotelBooking.ipynb}. \TBNotebook
\bibitem{tb-page} Group G3. \emph{Project page (rendered notebook, report and video)}. \TBPage\ — video: \TBVideo
\end{thebibliography}


\chapter{Text data: Vietnamese news topic classification}\label{ch:text}
% Body of the Text report — used by text.tex (as an article) and G3-report.tex (as a chapter).
% Every number comes from generated/text.tex (built from text/results.json).

\section{Problem and objective}\label{tx:problem}
News websites tag every article with a section so that readers can browse, search and receive recommendations.
Given only the \textbf{headline} and the \textbf{lead paragraph} (``sapo'') of a Vietnamese news article, we predict which of
\TXClasses{} newspaper sections it belongs to: News (Thời sự), World, Business, Sci-Tech, Entertainment, Sports, Law,
Education, Health and Travel.

\textbf{Objective.} Build a topic classifier for short Vietnamese texts, compare classical bag-of-words models with a
pretrained multilingual sentence-embedding model, and understand which mistakes remain.

\textbf{How success is measured.} \textbf{Macro-F1} — the mean of the per-section F1 scores, so every section counts
equally — plus accuracy, on a held-out test set used once.

\section{Dataset: collected by the group}\label{tx:dataset}
No ready-made dataset was used: the group wrote a crawler — \code{crawl\_vnexpress.py} in \code{text/crawler/}, built on
\code{requests} and \code{BeautifulSoup} — for \emph{VnExpress} \cite{tx-vnexpress}, one of Vietnam's largest online newspapers, and ran it on
27~September~2026.

\begin{figure}[H]\centering
\begin{tikzpicture}[node distance=0.36cm, box/.style={draw=bklight, fill=tint, rounded corners=2pt, font=\sffamily\footnotesize,
  text width=2.3cm, align=center, minimum height=1.25cm}, arr/.style={-{Stealth[length=2mm]}, bkdark, thick}]
\node[box] (a) {10 sections\\× 12 listing pages\\(120 requests)};
\node[box, right=of a] (b) {\code{requests}\\1 s pause between\\requests};
\node[box, right=of b] (c) {\code{BeautifulSoup}:\\headline, lead,\\dateline, URL};
\node[box, right=of c] (d) {raw CSV\\\TXRowsRaw{} rows\\(kept unmodified)};
\node[box, right=of d] (e) {notebook:\\clean, split,\\model};
\draw[arr] (a) -- (b); \draw[arr] (b) -- (c); \draw[arr] (c) -- (d); \draw[arr] (d) -- (e);
\end{tikzpicture}
\caption{Data collection pipeline. Only headlines, leads and URLs are stored (no article bodies), for non-commercial coursework.\label{tx:fig-pipeline}}
\end{figure}

\begin{table}[H]\centering\small
\caption{Fields of the crawled dataset.}\label{tx:tab-fields}
\begin{tabularx}{\linewidth}{L{3.6cm}Y}
\toprule \hd{Column} & \hd{Meaning} \\ \midrule
\code{title} & headline \\
\code{description} & lead paragraph (sapo) \\
\code{location} & dateline such as ``Hà Nội'' or ``Indonesia'', when present (not used as a feature) \\
\code{url} & article URL — unique key \\
\code{category\_slug}, \code{category} & section = \textbf{label} \\
\code{page}, \code{crawled\_at} & listing page number and crawl time \\
\bottomrule
\end{tabularx}
\end{table}

\begin{table}[H]\centering\small
\caption{Assignment requirements for text data (section 4.3) and how this dataset meets them.}\label{tx:tab-req}
\begin{tabularx}{\linewidth}{L{4.6cm}Yl}
\toprule \hd{Requirement} & \hd{This dataset} & \hd{Status} \\ \midrule
At least 2,000 samples & \TXRowsRaw{} articles (\TXRowsClean{} after cleaning) & \met \\
Vietnamese data encouraged & all headlines and leads are Vietnamese & \met \\
Self-collected data valued & crawled by the group; the crawler is in the repository & \met \\
\bottomrule
\end{tabularx}
\end{table}

\section{Exploratory data analysis}\label{tx:eda}
\textbf{Quality.} \TXMissingDesc{} articles have no lead paragraph (mostly photo or video stories); \TXDupWithin{} URLs appear
twice within the same section (an article was pinned while the listing shifted between requests); \TXMultiSection{} URLs
are listed under two different sections, which makes their label ambiguous. \TXNFC{} headlines needed Unicode normalisation.

\textbf{Balance and length.} The crawler read the same number of pages per section, so the classes are balanced
(\TXClassMin{}–\TXClassMax{} articles each). Texts are short: a median headline has \TXMedTitle{} syllables and a median
lead \TXMedDesc{} — about the length of one tweet. Vietnamese separates \emph{syllables}, not words, with spaces
(``bóng đá'', football, is two tokens), which shapes the choice of representation (Section~\ref{tx:prep}).

\begin{table}[H]\centering\small
\caption{The eight n-grams most over-represented in each section (χ² test of each TF-IDF 1–2-gram, section vs.\ rest, training set).\label{tx:tab-terms}}
% Generated by tools/build_reports.py — do not edit.
\begin{tabularx}{\linewidth}{lY}
\toprule
\hd{Section} & \hd{Top 8 discriminative n-grams (χ², training set)} \\
\midrule
Business & lãi, phiếu, doanh nghiệp, cổ phiếu, lãi suất, doanh, khoán, chứng khoán \\
Education & học, đại học, sinh, trường, học sinh, sinh viên, đại, ngành \\
Entertainment & phim, miss, diễn viên, miss world, diễn, hoa hậu, ca, hậu \\
Health & bệnh, huyết, bác sĩ, máu, viêm, ăn, uống, ung \\
Law & bị cáo, cáo, bị, cáo buộc, công an, sát hại, tù, lừa \\
News & ngập, mưa, mưa lớn, sông, quốc lộ, lộ num, sạt, lở \\
Sci-Tech & iphone, robot, num pro, pro, iphone num, apple, ai, robot hình \\
Sports & cup, hlv, asiad, man, trận, asiad num, league, fifa \\
Travel & khách, du, du khách, du lịch, điểm đến, khách việt, lịch, dịp \\
World & trump, iran, ukraine, ông trump, tổng thống, tổng, nga, houthi \\
\bottomrule
\end{tabularx}

\end{table}
\begin{finding}
\Q Which words does each section ``own''?
\F Two kinds of signal appear. \textbf{Stable topic vocabulary}: \emph{Health} — \textit{bệnh, bác sĩ, viêm}; \emph{Education}
— \textit{học sinh, đại học, sinh viên}; \emph{Law} — \textit{bị cáo, cáo buộc, tù}; \emph{Business} — \textit{cổ phiếu, lãi
suất, chứng khoán}. Bigrams such as \textit{cổ phiếu} and \textit{bác sĩ} rank next to their single syllables, so syllable
pairs recover real words. \textbf{Time-specific names}: \emph{Sci-Tech} is driven by \textit{iphone, apple, robot},
\emph{World} by \textit{trump, iran, ukraine}, \emph{Sports} by \textit{asiad, fifa}, and \emph{News} by recent floods
(\textit{ngập, mưa lớn, sạt lở}). These reflect the news cycle at crawl time rather than the section itself — a risk
discussed in the limitations.
\end{finding}

\fig[0.72\linewidth]{txt_03_centroid_similarity}{Cosine similarity between the average TF-IDF vector (centroid) of each section.\label{tx:fig-similarity}}
\begin{finding}
\Q Before training any model, can we predict which sections will be confused?
\F The most similar pairs are \textbf{\TXSimOnePair{}} (\TXSimOne) — chips, technology companies and market value appear in
both — and \textbf{\TXSimTwoPair{}} (\TXSimTwo) — police and accidents appear in both. \emph{Sports} is the most isolated
section. Section~\ref{tx:results} checks this prediction against the confusion matrix.
\end{finding}

\section{Text preparation}\label{tx:prep}
\begin{table}[H]\centering\small
\caption{Preparation steps.}\label{tx:tab-prep}
\begin{tabularx}{\linewidth}{L{4.0cm}Y}
\toprule \hd{Step} & \hd{What is done and why} \\ \midrule
Duplicates & keep one copy of URLs repeated within a section; drop URLs listed under two sections (ambiguous label) → \TXRowsClean{} articles (\TXRemoved{} rows removed) \\
Missing values & a missing lead becomes an empty string, so the article is classified from its headline alone \\
Input text & headline + ``. '' + lead \\
Cleaning (\code{VietnameseTextCleaner}) & Unicode NFC normalisation (the same diacritic can be stored as one or two code points); lower-casing; numbers → token \code{num}; punctuation removed \\
Tokens & pattern \code{\textbackslash b\textbackslash w+\textbackslash b} keeps one-character syllables; the scikit-learn default would silently drop \TXSingleChar{} of them, including \textit{ô} (in \textit{ô tô}, car) and \textit{ở} (at) \\
Label encoding & \code{LabelEncoder}: 10 sections → 0–9 \\
Split & stratified 70 / 15 / 15: \TXTrain{} / \TXVal{} / \TXTest{} articles \\
\bottomrule
\end{tabularx}
\end{table}

For text, \emph{encoding} means turning each document into a vector and \emph{scaling} means weighting and normalising that
vector. Four representations are compared, all fitted on the training set only:

\begin{table}[H]\centering\small
\caption{Text representations compared.}\label{tx:tab-reps}
\begin{tabularx}{\linewidth}{L{3.0cm}YL{3.3cm}r}
\toprule \hd{Representation} & \hd{Idea} & \hd{Weighting / scaling} & \hd{Vocabulary} \\ \midrule
BoW, 1-gram & how often each syllable occurs & raw counts & \TXVocabUni \\
TF-IDF, 1-gram & counts re-weighted by rarity across documents & sublinear tf, L2 norm & \TXVocabUni \\
TF-IDF, 1–2-gram & adds syllable pairs → compound words & sublinear tf, L2 norm & \TXVocabBi \\
TF-IDF, char 2–5 & character n-grams inside words → robust to spelling & sublinear tf, L2 norm & \TXVocabChar \\
\bottomrule
\end{tabularx}
\end{table}

\section{Modelling}\label{tx:models}
\subsection{Representation × classifier grid}
Each representation is paired with four classifiers — Multinomial Naive Bayes, Logistic Regression, a linear Support Vector
Machine and k-nearest neighbours with cosine distance — and scored on the validation set (\code{TextBenchmark} class).

\begin{table}[H]\centering\small
\caption{Validation macro-F1 for every representation × classifier pair (best in bold). A stratified random guess scores \TXRandomVal.}\label{tx:tab-grid}
% Generated by tools/build_reports.py — do not edit.
\begin{tabular}{lrrrr}
\toprule
\hd{Representation} & \hd{Multinomial NB} & \hd{Logistic Reg.} & \hd{Linear SVM} & \hd{kNN (cosine)} \\
\midrule
BoW 1-gram & 0.839 & 0.820 & 0.819 & 0.751 \\
TF-IDF 1-gram & 0.837 & 0.848 & 0.858 & 0.837 \\
TF-IDF 1–2-gram & 0.856 & \textbf{0.864} & 0.863 & 0.849 \\
TF-IDF char 2–5 & 0.825 & 0.847 & 0.858 & 0.819 \\
\bottomrule
\end{tabular}

\end{table}
\begin{finding}
\Q Which representation and classifier work best for short Vietnamese texts?
\F \textbf{TF-IDF beats raw counts} for the discriminative models — Logistic Regression \TXGainTfidfLr, Linear SVM
\TXGainTfidfSvm, kNN \TXGainTfidfKnn{} — because IDF down-weights syllables that occur everywhere (\textit{của, là, và}),
acting as an automatic stop-word list; Naive Bayes, designed for counts, does not gain. \textbf{Adding bigrams helps every
model} (\TXGainBiMin{} to \TXGainBiMax), while character n-grams do not beat word bigrams. The best pair is
\TXGridBest{} (\TXGridBestF); the Linear SVM on the same features is practically tied (\TXGBiSvm) and trains several times faster.
\end{finding}

\subsection{Tuning}
The TF-IDF + Linear SVM pipeline was tuned with 5-fold stratified cross-validation on the training set over the n-gram range,
\code{min\_df}, sublinear term frequency and the regularisation constant \code{C}. The best setting uses
\TXBestNgram-grams, \code{min\_df} = \TXBestMinDf{} and \code{C} = \TXBestC{} (cross-validated macro-F1 \TXCvF). On the single
validation split it scores \TXTunedVal{} — within about 0.01 of the untuned models, i.e.\ a handful of the \TXVal{} articles;
the 5-fold estimate is the more reliable of the two.

\section{Extra experiments}\label{tx:extra}
\subsection{Which field carries the signal, and would more data help?}
\begin{figure}[H]\centering
\begin{subfigure}[t]{0.49\linewidth}\centering\includegraphics[width=\linewidth]{txt_06_fields}
\caption{Headline vs lead as input (validation).}\end{subfigure}\hfill
\begin{subfigure}[t]{0.49\linewidth}\centering\includegraphics[width=\linewidth]{txt_07_learning_curve}
\caption{Learning curve of the tuned model.}\end{subfigure}
\caption{Two questions about the data itself, answered with the tuned TF-IDF + Linear SVM pipeline.\label{tx:fig-data}}
\end{figure}
\begin{finding}
\Q Is the headline or the lead more informative — and is our crawled dataset big enough?
\F The \textbf{lead alone (\TXLeadOnly) beats the headline alone (\TXHeadOnly)} by \TXLeadOverHead: leads are longer and
written to summarise, headlines are short and often figurative. Using both is best (\TXBoth, \TXBothOverLead{} over the lead).
The learning curve rises from \TXLcFirst{} with \TXLcFirstN{} articles to \TXLcMid{} with \TXLcMidN, then flattens: doubling to
\TXLcLastN{} articles adds only \TXLcLastGain. Crawling more would help a little; a better representation should help more.
\end{finding}

\subsection{TF-IDF versus multilingual sentence embeddings}
We encode every article with \code{intfloat/multilingual-e5-small} \cite{tx-e5}, a 118-million-parameter multilingual sentence
encoder with 384-dimensional output, loaded with the sentence-transformers library \cite{tx-sbert}. Following the assignment, the
pretrained model is used as a \textbf{frozen feature extractor}: a Logistic Regression is trained on the embeddings. A
\textbf{hybrid} model concatenates the TF-IDF features with the embeddings and trains a Linear SVM.

\fig[0.8\linewidth]{txt_08_tfidf_vs_embedding}{Test macro-F1 of the three approaches and the random baseline (dark = final model, chosen on validation).\label{tx:fig-embed}}
\begin{finding}
\Q Does a pretrained neural encoder beat a bag of syllables on short Vietnamese news?
\F Yes, and combining them is best. On the test set, embeddings alone reach \TXEmbTest{} (\TXGainEmb{} over tuned TF-IDF at
\TXTfidfTest), and the \textbf{hybrid reaches \TXHybTest{}} (\TXGainHyb). The ranking is the same on validation
(\TXTfidfVal{} → \TXEmbVal{} → \TXHybVal), so the choice generalises. TF-IDF keeps exact names (\textit{iphone, asiad}) that
embeddings blur; embeddings map different wordings of the same idea close together.
\end{finding}

\section{Results on the test set}\label{tx:results}
\begin{table}[H]\centering\small
\caption{Validation and test scores of the compared approaches. The final model was chosen on validation macro-F1.}\label{tx:tab-comparison}
% Generated by tools/build_reports.py — do not edit.
\begin{tabular}{lrrr}
\toprule
\hd{Approach} & \hd{Val. macro-F1} & \hd{Test macro-F1} & \hd{Test acc.} \\
\midrule
Random guess (stratified) & — & 0.102 & — \\
TF-IDF 1–3-gram + Linear SVM (tuned) & 0.857 & 0.865 & 0.866 \\
Embedding (e5-small) + Logistic Regr. & 0.882 & 0.873 & 0.873 \\
\textbf{Hybrid: TF-IDF ⊕ embedding + Linear SVM} & \textbf{0.888} & \textbf{0.892} & \textbf{0.892} \\
\bottomrule
\end{tabular}

\end{table}

The final hybrid model reaches \textbf{accuracy \TXFinAcc{} and macro-F1 \TXFinF} on the \TXTest{} test articles; \TXNMis{} are
misclassified. \emph{\TXBestClass} is the easiest section (F1 \TXBestClassF) and \emph{\TXWorstClass} the hardest
(\TXWorstClassF).

\begin{table}[H]\centering\small
\caption{Per-section test F1 of the final model.}\label{tx:tab-perclass}
% Generated by tools/build_reports.py — do not edit.
\begin{tabular}{lrlr}
\toprule
\hd{Section} & \hd{F1} & \hd{Section} & \hd{F1} \\
\midrule
Sports & 0.973 & Travel & 0.878 \\
Entertainment & 0.946 & News & 0.865 \\
Health & 0.937 & World & 0.846 \\
Education & 0.924 & Sci-Tech & 0.844 \\
Law & 0.893 & Business & 0.814 \\
\bottomrule
\end{tabular}

\end{table}

\fig[0.7\linewidth]{txt_09_confusion}{Confusion matrix of the final model on the test set (counts; colour = share of the true section).\label{tx:fig-confusion}}

\begin{table}[H]\centering\small
\caption{Most frequent confusions on the test set.}\label{tx:tab-confusions}
% Generated by tools/build_reports.py — do not edit.
\begin{tabular}{llr}
\toprule
\hd{True section} & \hd{Predicted as} & \hd{Test errors} \\
\midrule
News & Education & 4 \\
Business & Sci-Tech & 4 \\
World & Law & 3 \\
Sci-Tech & Travel & 3 \\
World & Business & 3 \\
\bottomrule
\end{tabular}

\end{table}
\begin{finding}
\Q Did the centroid similarity of Figure~\ref{tx:fig-similarity} predict the errors?
\F Largely yes: \textbf{Business → Sci-Tech}, the most similar pair before training, is among the two most frequent confusions.
Business is the hardest section because it overlaps with technology, world affairs and travel (airlines, tourism revenue).
\end{finding}

\begin{table}[H]\centering\small
\caption{Examples of misclassified test articles.}\label{tx:tab-wrong}
% Generated by tools/build_reports.py — do not edit.
\begin{tabularx}{\linewidth}{Yll}
\toprule
\hd{Headline (Vietnamese)} & \hd{True} & \hd{Predicted} \\
\midrule
Thủ đô quốc gia Đông Nam Á nào đông dân nhất thế giới? & Travel & Education \\
Giải cứu 10 ngư dân Việt Nam làm việc tại Malaysia & World & News \\
Hàng không đón khách quốc tế nhiều hơn nội địa dịp Quốc khánh & Business & Travel \\
Tại sao cửa sổ máy bay khó sản xuất? & Sci-Tech & Travel \\
Nghệ thuật 'khiêu khích' ẩn sau quả chuối dán tường 6 triệu USD & World & Travel \\
Nỗi day dứt của đầu bếp thoát chết vụ 11/9 nhờ đổi ca với đồng nghiệp & World & Law \\
Justin Bieber mở show miễn phí vì người vô gia cư & Entertainment & Travel \\
Mark Zuckerberg mua lâu đài & Sci-Tech & Travel \\
\bottomrule
\end{tabularx}

\end{table}
Many errors sit genuinely between sections — \textit{``Tại sao cửa sổ máy bay khó sản xuất?''} (why aircraft windows are hard to
make) is Sci-Tech but reads like Travel; a human editor could file several of them either way, so part of the remaining error
is label ambiguity rather than model failure. Finally, all eight headlines written by the group, which are not in the dataset,
are classified correctly:

\begin{table}[H]\centering\small
\caption{Predictions of the final model on new headlines.}\label{tx:tab-demo}
% Generated by tools/build_reports.py — do not edit.
\begin{tabularx}{\linewidth}{Yl}
\toprule
\hd{Headline written by the group (not in the dataset)} & \hd{Predicted section} \\
\midrule
Đội tuyển Việt Nam thắng Thái Lan 2-1 ở chung kết & Sports \\
Giá vàng tăng mạnh phiên đầu tuần & Business \\
Bộ Giáo dục công bố phương án thi tốt nghiệp THPT 2027 & Education \\
Phát hiện hành tinh mới có thể chứa nước & Sci-Tech \\
Bác sĩ khuyến cáo cách phòng sốt xuất huyết mùa mưa & Health \\
Khởi tố giám đốc công ty lừa đảo chiếm đoạt 50 tỷ đồng & Law \\
Tổng thống Mỹ gặp Chủ tịch Trung Quốc tại hội nghị APEC & World \\
Vịnh Hạ Long lọt top điểm đến đẹp nhất châu Á & Travel \\
\bottomrule
\end{tabularx}

\end{table}

\section{Conclusions}\label{tx:conclusion}
\begin{enumerate}
\item \textbf{A self-crawled Vietnamese dataset is enough for a strong classifier}: \TXRowsClean{} clean articles in \TXClasses{}
sections give \TXFinF{} test macro-F1 from a headline and a two-sentence lead.
\item \textbf{Vietnamese needs care in the details}: NFC normalisation, keeping one-character syllables, and syllable bigrams to
recover compound words. TF-IDF 1–2-grams beat raw unigram counts by \TXGainBiOverBowLr{} (Logistic Regression) and
\TXGainBiOverBowSvm{} (Linear SVM).
\item \textbf{Sentence embeddings beat TF-IDF, and the combination is best}: \TXGainEmb{} and \TXGainHyb{} test macro-F1 over tuned TF-IDF.
\item \textbf{The lead carries more signal than the headline} (\TXLeadOnly{} vs \TXHeadOnly), and the learning curve flattens
after about \TXLcMidN{} articles.
\item \textbf{Errors are predictable}: the most similar pair of section centroids is among the most frequent confusions.
\end{enumerate}
\textbf{Limitations.} All articles come from one short window (the latest 12 listing pages per section on 27~September~2026),
and many discriminative words are event names (Trump, ASIAD, iPhone, floods), so accuracy will likely drop on news from
another period. There is one newspaper and one editorial style, and section labels are sometimes ambiguous. Syllable
n-grams only approximate words; a Vietnamese word segmenter was not used, to keep the notebook dependency-free.

\textbf{Future work.} Re-crawl in a later month and measure the drop (temporal robustness); crawl other newspapers to test
cross-source generalisation; compare word-segmented TF-IDF; lightly fine-tune a Vietnamese language model on top of the
frozen-encoder baseline.

\begin{thebibliography}{9}
\bibitem{tx-vnexpress} VnExpress — Vietnamese online newspaper. \url{https://vnexpress.net} (section listings crawled 27~September~2026).
\bibitem{tx-e5} L.~Wang, N.~Yang, X.~Huang, L.~Yang, R.~Majumder, F.~Wei. \emph{Multilingual E5 Text Embeddings: A Technical
  Report}. arXiv:2402.05672, 2024. Model: \url{https://huggingface.co/intfloat/multilingual-e5-small}
\bibitem{tx-sbert} N.~Reimers, I.~Gurevych. \emph{Sentence-BERT: Sentence Embeddings using Siamese BERT-Networks}. EMNLP 2019.
\bibitem{tx-sklearn} F.~Pedregosa et al. \emph{Scikit-learn: Machine Learning in Python}. Journal of Machine Learning
  Research, 12 (2011), 2825--2830.
\bibitem{tx-nb} Group G3. \emph{Notebook} \code{G3\_Text\_VnExpressNews.ipynb}. \TXNotebook
\bibitem{tx-page} Group G3. \emph{Project page (rendered notebook, report and video)}. \TXPage\ — video: \TXVideo
\end{thebibliography}


\chapter{Time-series data: forecasting Hanoi's daily temperature}\label{ch:timeseries}
% Body of the Time-series report — used by timeseries.tex (as an article) and G3-report.tex (as a chapter).
% Every number comes from generated/timeseries.tex (built from timeseries/results.json).

\section{Problem and objective}\label{ts:problem}
Hanoi has a real winter: cold surges of the north-east monsoon can cut the temperature by 5–10~°C within a day, and those
are exactly the days when a forecast matters most for health, energy use and agriculture.

\textbf{Objective.} Forecast \textbf{tomorrow's daily mean temperature in Hanoi} from today's and past weather, then measure
how quickly forecast skill decays when predicting up to seven days ahead.

\textbf{How success is measured.} Mean absolute error (MAE, °C) on a chronologically held-out test period, and \emph{skill}
relative to simple baselines — above all \emph{persistence} (``tomorrow = today''), which is hard to beat for one-day forecasts.

\section{Dataset}\label{ts:dataset}
The data was downloaded by the group from the \textbf{Open-Meteo Historical Weather API} \cite{ts-openmeteo} (free, no API
key, CC BY 4.0), which serves the \textbf{ERA5 reanalysis} \cite{ts-era5}: a physically consistent reconstruction of past
weather on a 9–25~km grid, combining a weather model with millions of observations. The exact request is included in the
notebook (\code{fetch\_open\_meteo}); \emph{Run all} uses the saved CSV so results are reproducible.

\begin{itemize}
\item Location: central Hanoi (21.03~°N, 105.85~°E; nearest grid cell 21.05~°N, 105.90~°E, elevation 19~m).
\item Period: \TSStart{} – \TSEnd{}, daily, local time (GMT+7): \TSDays{} consecutive days, \TSVars{} variables.
\end{itemize}

\begin{table}[H]\centering\small
\caption{Variables. The target is the daily mean temperature.}\label{ts:tab-vars}
\begin{tabularx}{\linewidth}{L{6.4cm}lY}
\toprule \hd{Variable} & \hd{Unit} & \hd{Meaning} \\ \midrule
\code{temperature\_2m\_mean} / \code{\_max} / \code{\_min} & °C & daily mean (\textbf{target}), maximum, minimum at 2~m \\
\code{apparent\_temperature\_mean} & °C & ``feels-like'' temperature \\
\code{precipitation\_sum}, \code{precipitation\_hours} & mm, h & daily rain and rainy hours \\
\code{relative\_humidity\_2m\_mean}, \code{dew\_point\_2m\_mean} & \%, °C & moisture \\
\code{pressure\_msl\_mean} & hPa & mean sea-level pressure \\
\code{cloud\_cover\_mean} & \% & cloud cover \\
\code{wind\_speed\_10m\_max}, \code{wind\_direction\_10m\_dominant} & km/h, ° & maximum wind speed, dominant direction \\
\code{shortwave\_radiation\_sum}, \code{sunshine\_duration} & MJ/m², s & solar energy, sunshine \\
\bottomrule
\end{tabularx}
\end{table}

\begin{table}[H]\centering\small
\caption{General requirements (section 4.1 and 5) and how this dataset meets them.}\label{ts:tab-req}
\begin{tabularx}{\linewidth}{L{4.2cm}Yl}
\toprule \hd{Requirement} & \hd{This dataset} & \hd{Status} \\ \midrule
Enough samples & \TSDays{} days & \met \\
Numerical columns (scaling) & \TSVars{} weather variables plus engineered lags & \met \\
Categorical columns (encoding) & derived: season and 8-sector wind direction (one-hot), circular encodings & \met \\
Missing values handled & none in the source (verified); gaps simulated to compare imputation methods (Section~\ref{ts:missing}) & \met \\
Time-aware split & train / validation / test by date; time-series cross-validation & \met \\
\bottomrule
\end{tabularx}
\end{table}

\section{Exploratory data analysis}\label{ts:eda}
\fig{ts_02_seasonality}{Left: climatology by day of year (mean, ±1 standard deviation, record range). Right: mean absolute day-to-day change by month.\label{ts:fig-season}}
\begin{finding}
\Q How different are the seasons — and when is the weather hardest to predict?
\F Summer (June–August) averages \textbf{\TSSummer~°C} and winter (December–February) \textbf{\TSWinter~°C}, a swing of about
\TSSeasonGap~°C. Winter is also far more \emph{variable}: the standard deviation within January is \TSJanStd~°C versus
\TSJulStd~°C in July, and the average day-to-day change peaks in February–March (\TSDtdFeb–\TSDtdMar~°C) against \TSDtdJul~°C
in July. Forecasting should therefore be hardest in winter and spring.
\end{finding}

\fig{ts_03_trend}{Left: annual mean temperature with a linear trend. Right: monthly anomaly versus the 2005–2024 average (blue = colder, red = warmer).\label{ts:fig-trend}}
\begin{finding}
\Q Is Hanoi getting warmer?
\F The annual mean rises by \textbf{\TSTrend~°C per decade} over 2005–2025; 2021–25 averaged \TSLastFive~°C against
\TSFirstFive~°C in 2005–09 (\TSDeltaFive~°C). Large cold anomalies cluster in the early years (February 2008, January–February
2011) and recent years are mostly warm. Twenty-one years is short and year-to-year swings are large (\TSAnnualMin~°C in
\TSAnnualMinYear{} vs \TSAnnualMax~°C in \TSAnnualMaxYear), so the slope is uncertain — but it matters for modelling: a
climatology learned on 2005–2021 will be slightly too cold for 2024–2026 (Section~\ref{ts:results}).
\end{finding}

\begin{figure}[H]\centering
\begin{subfigure}[t]{0.58\linewidth}\centering\includegraphics[width=\linewidth]{ts_04_acf}
\caption{Autocorrelation of temperature and of its anomaly.}\end{subfigure}\hfill
\begin{subfigure}[t]{0.4\linewidth}\centering\includegraphics[width=\linewidth]{ts_05_cold_surge}
\caption{Pressure change today vs temperature change tomorrow (Nov–Mar).}\end{subfigure}
\caption{Memory of the series and the cold-surge signal.\label{ts:fig-memory}}
\end{figure}
\begin{finding}
\Q How much does today tell us about tomorrow, and can we see a cold surge coming?
\F The raw series is dominated by the yearly cycle (autocorrelation \TSAcfYear{} at lag 365). Once the season is removed, the
anomaly still has strong day-to-day memory (\textbf{\TSAcfAnomOne{} at lag 1}) that fades to \textbf{\TSAcfAnomSeven{} by lag 7}:
a warm or cold spell ``remembers'' itself for only a few days. In the cold season a pressure rise today warns of cooling
tomorrow (r = \TSR): on the \TSNRise{} days when pressure rose by more than 3~hPa, the next day changed by \TSDTRise~°C on
average (vs \TSDTOther~°C otherwise) and a drop of more than 2~°C was \textbf{\TSPRatio× as likely} (\TSPRise{} vs \TSPOther).
\end{finding}

\section{Data preparation}\label{ts:prep}
\subsection{Continuity and missing values}\label{ts:missing}
A time series must be checked for missing \emph{dates}, not only missing cells, because a skipped day silently shifts every
lag feature. The file has \TSDays{} rows for \TSDays{} calendar days, no duplicated dates and \TSMissingDates{} missing cells.
Real station data is rarely this clean, so we used the complete series as ground truth for an experiment: hide random blocks of
1–30 consecutive days (about 5\% of the series, 5 repetitions) and reconstruct them with four methods.

\begin{table}[H]\centering\small
\caption{Imputation error (MAE in °C on the hidden days; best per gap length in bold).}\label{ts:tab-imputation}
% Generated by tools/build_reports.py — do not edit.
\begin{tabular}{lrrrr}
\toprule
\hd{Gap length} & \hd{Forward fill} & \hd{Linear interp.} & \hd{Climatology} & \hd{Seasonal interp.} \\
\midrule
1 day & 1.16 & 0.71 & 1.97 & \textbf{0.70} \\
3 days & 1.76 & 1.14 & 1.89 & \textbf{1.13} \\
7 days & 2.21 & 1.69 & 2.00 & \textbf{1.63} \\
14 days & 2.45 & 1.95 & \textbf{1.86} & 1.93 \\
30 days & 2.80 & 2.22 & \textbf{1.87} & 2.15 \\
\bottomrule
\end{tabular}

\end{table}
\begin{finding}
\Q If a station lost data for a few days, how should the gap be filled — and does the answer depend on its length?
\F For short gaps, \textbf{seasonal interpolation} (interpolating the anomaly, then adding the climatology back) is best:
\TSImpSeasOne, \TSImpSeasThree{} and \TSImpSeasSeven~°C for 1, 3 and 7 days, slightly ahead of plain linear interpolation
(\TSImpLinOne, \TSImpLinThree, \TSImpLinSeven). Forward fill is worst at every length. For gaps of \textbf{14 days or more,
the seasonal normal alone wins} (\TSImpClimFourteen~°C), because the anomaly's memory (about a week) has run out. Rule:
interpolate the anomaly up to about a week, fall back to climatology for longer outages.
\end{finding}

\subsection{Outliers}
Temperature outliers must be judged against the season (12~°C is normal in January but extreme in July), so each day's anomaly
is divided by the standard deviation for that day of the year and |z| > 3 is flagged. Only \TSExtremes{} days qualify, all
cold (March 2011, early October 2011, September 2013): real cold spells, exactly what a forecast must handle, so they are
\textbf{kept}. The record \TSColdest~°C of \TSColdestDate{} is not even flagged because January is so variable.
Precipitation has \TSRainOutliers{} IQR outliers and skewness \TSRainSkew{} (maximum \TSRainMax~mm on \TSRainMaxDate); heavy
rain is also real, so it is kept and fed to the models as \code{log1p(rain)} (skewness \TSRainSkewLog).

\subsection{Features, encoding and scaling}
For a forecast issued at the end of day \emph{t} for day \emph{t + h}, every feature must be known on day \emph{t}. The class
\code{WeatherFeatureBuilder} builds them; its climatology is learned from the \textbf{training years only}.

\begin{table}[H]\centering\small
\caption{Features (30 in total) and their encoding.}\label{ts:tab-features}
\begin{tabularx}{\linewidth}{L{3.5cm}YL{3.4cm}}
\toprule \hd{Group} & \hd{Features} & \hd{Encoding / scaling} \\ \midrule
Recent temperature & today; lags 1, 2, 3, 7, 14 days; rolling means over 3, 7, 30 days; today's max, min and range & standardised \\
Anomaly & today − climatology (smoothed day-of-year mean) & standardised \\
Weather today & humidity, dew point, pressure, pressure change over 1 and 3 days, cloud cover, wind speed, rain, radiation, sunshine & \code{log1p} for rain; standardised \\
Wind direction & sine and cosine of the direction & circular (359° is next to 1°) \\
Calendar of target day & day of year; climatology of the target day & sine/cosine (cyclical) \\
Categorical & season of target day (DJF, MAM, JJA, SON); 8-sector wind direction (N, NE, …, NW) & one-hot \\
\bottomrule
\end{tabularx}
\end{table}
The first 30 days (rolling-window warm-up) and the last day (no target) are dropped. Scaling uses \code{StandardScaler} inside
the pipeline; it is required by Ridge regression and harmless for tree models.

\subsection{Time-based split}
A random split would let the model learn from days \emph{after} the ones it is tested on. We split by the date being forecast:
\textbf{train} 2005–2021 (\TSTrain{} days), \textbf{validation} 2022–2023 (\TSVal{} days) and \textbf{test} 1~January~2024 –
\TSEnd{} (\TSTest{} days). Hyperparameters are tuned with \code{TimeSeriesSplit} (5 expanding windows inside the training years),
which always trains on the past and validates on the following period.

\section{Modelling}\label{ts:models}
\textbf{Baselines.} \emph{Persistence} (tomorrow = today); \emph{climatology} (the average for that calendar day); and
\emph{damped anomaly persistence} (climatology + φ × today's anomaly, with φ = \TSPhi, the lag-1 autocorrelation of anomalies
in the training years).

\textbf{Machine-learning models.} Ridge regression, Random Forest and HistGradientBoosting (HGB), all in the same
pipeline. HGB is trained either on the temperature itself or on the \textbf{change} ΔT = T(t+1) − T(t), adding today's value
back afterwards (\code{DeltaRegressor} class). Trees cannot extrapolate beyond the range seen in training, and a warming climate
makes that a real risk; predicting the change avoids it.

\textbf{Tuning.} A random search over 20 HGB configurations (scikit-learn defaults included) with time-series cross-validation
preferred a more regularised model (CV MAE \TSCvTuned{} vs \TSCvDefault{} for the defaults), but it did \emph{worse} on
2022–2023 (\TSValTuned{} vs \TSValFinal). A 0.015~°C gain on older folds is within year-to-year noise, so the final model is
the one chosen on validation: HGB on the change, with default hyperparameters.

\begin{table}[H]\centering\small
\caption{One-day-ahead forecasts: validation (2022–2023) and test (2024 – September 2026) scores. Skill = 1 − MAE / MAE of persistence (test).}\label{ts:tab-models}
% Generated by tools/build_reports.py — do not edit.
\begin{tabular}{lrrrrr}
\toprule
\hd{Model} & \hd{Val MAE} & \hd{Test MAE} & \hd{Test RMSE} & \hd{Test R²} & \hd{Skill} \\
\midrule
Climatology (seasonal normal) & 2.014 & 1.990 & 2.538 & 0.722 & −84.2\% \\
Persistence (tomorrow = today) & 1.105 & 1.081 & 1.504 & 0.902 & 0.0\% \\
Damped anomaly persistence & 1.092 & 1.075 & 1.434 & 0.911 & 0.5\% \\
Ridge regression & 0.974 & 0.970 & 1.283 & 0.929 & 10.2\% \\
Random Forest & 0.921 & 0.935 & 1.245 & 0.933 & 13.4\% \\
HGB, temperature as target & 0.933 & 0.932 & 1.231 & 0.934 & 13.8\% \\
HGB, change as target, tuned & 0.922 & 0.919 & 1.222 & 0.935 & 15.0\% \\
\textbf{HGB, change as target (final)} & \textbf{0.906} & \textbf{0.916} & \textbf{1.226} & \textbf{0.935} & \textbf{15.3\%} \\
\bottomrule
\end{tabular}

\end{table}
\begin{finding}
\Q Do machine-learning models beat the classical baselines, and which framing works best?
\F Every model beats every baseline. HGB predicting the change is best on validation (\TSValFinal~°C, \TSValSkill{} below
persistence at \TSValPersist) and beats HGB predicting the level (\TSValLevel). Ridge regression is weakest (\TSValRidge): the
useful signals interact non-linearly — a northerly wind matters mostly in winter. Damped anomaly persistence barely improves on
persistence (\TSValDamped) and climatology alone is poor (\TSValClim~°C): for \emph{tomorrow}, today's weather matters far more
than the calendar.
\end{finding}

\section{Results on the test set}\label{ts:results}
On the \TSTest{} unseen days the final model reaches \textbf{MAE \TSTestFinal~°C, RMSE \TSTestRmse~°C and R² \TSTestRtwo},
a \textbf{\TSTestSkill{} lower error than persistence} (\TSTestPersist~°C); the ranking of all models matches validation
(Table~\ref{ts:tab-models}).

\fig[0.86\linewidth]{ts_10_forecast_vs_actual}{Forecast versus actual over the whole test period (top) and the hardest winter window (bottom), with persistence for comparison.\label{ts:fig-forecast}}

\begin{table}[H]\centering\small
\caption{Test MAE by month (bold = final model better than persistence).}\label{ts:tab-months}
\footnotesize\setlength{\tabcolsep}{3.6pt}
% Generated by tools/build_reports.py — do not edit.
\begin{tabular}{lrrrrrrrrrrrr}
\toprule
\hd{Test MAE (°C)} & \hd{Jan} & \hd{Feb} & \hd{Mar} & \hd{Apr} & \hd{May} & \hd{Jun} & \hd{Jul} & \hd{Aug} & \hd{Sep} & \hd{Oct} & \hd{Nov} & \hd{Dec} \\
\midrule
Persistence & 1.36 & 1.61 & 1.24 & 1.08 & 1.08 & 0.99 & 0.75 & 0.87 & 0.86 & 0.96 & 1.01 & 1.12 \\
Final model & \textbf{1.05} & \textbf{1.16} & \textbf{0.82} & 1.10 & \textbf{0.91} & \textbf{0.87} & \textbf{0.69} & \textbf{0.84} & \textbf{0.72} & \textbf{0.90} & 1.12 & \textbf{0.88} \\
\bottomrule
\end{tabular}

\end{table}

\begin{table}[H]\centering\small
\caption{The \TSWorstN{} largest test errors (°C; pressure change on the forecast day in hPa).}\label{ts:tab-worst}
% Generated by tools/build_reports.py — do not edit.
\begin{tabular}{lrrrrr}
\toprule
\hd{Date} & \hd{Actual} & \hd{Forecast} & \hd{Error} & \hd{1-day change} & \hd{Pressure change} \\
\midrule
29 March 2025 & 17.0 & 23.7 & +6.7 & −8.5 & +4.2 \\
8 February 2026 & 15.0 & 20.7 & +5.7 & −7.0 & +2.2 \\
18 November 2025 & 17.8 & 22.7 & +4.9 & −4.9 & −0.1 \\
23 February 2024 & 19.3 & 23.8 & +4.5 & −6.0 & +3.7 \\
26 January 2025 & 14.9 & 19.4 & +4.5 & −6.4 & +1.6 \\
29 April 2024 & 31.7 & 27.4 & −4.3 & +1.0 & +2.3 \\
7 February 2024 & 15.7 & 20.0 & +4.3 & −5.7 & +0.8 \\
22 January 2024 & 11.3 & 15.5 & +4.2 & −4.6 & +3.5 \\
\bottomrule
\end{tabular}

\end{table}
\begin{finding}
\Q When does the model fail?
\F The biggest gains over persistence come in winter and spring — February \TSMoFebModel{} vs \TSMoFebPersist, March
\TSMoMarModel{} vs \TSMoMarPersist, December \TSMoDecModel{} vs \TSMoDecPersist~°C — while the model is slightly worse in the
transition months (April \TSMoAprModel{} vs \TSMoAprPersist, November \TSMoNovModel{} vs \TSMoNovPersist). \textbf{\TSWorstCold{} of
the \TSWorstN{} worst forecasts are cold surges}: the temperature fell \TSWorstDropMin–\TSWorstDropMax~°C in one day and the
forecast was \TSWorstErrMin–\TSWorstErrMax~°C too warm, usually after a pressure rise that the model reacted to only partly. On
average the residual is \TSResMean~°C (standard deviation \TSResStd): forecasts are slightly too cold, consistent with the
warming of Figure~\ref{ts:fig-trend}.
\end{finding}

\section{Extra experiments}\label{ts:extra}
\subsection{How far ahead is the model useful?}
The features are rebuilt and the final configuration retrained for horizons of 1 to 7 days; all methods are scored on the same
test period.

\begin{table}[H]\centering\small
\caption{Test MAE (°C) by forecast horizon (best per horizon in bold).}\label{ts:tab-horizon}
% Generated by tools/build_reports.py — do not edit.
\begin{tabular}{lrrrrrrr}
\toprule
\hd{Test MAE (°C)} & \hd{1 d} & \hd{2 d} & \hd{3 d} & \hd{4 d} & \hd{5 d} & \hd{6 d} & \hd{7 d} \\
\midrule
Persistence (tomorrow = today) & 1.08 & 1.74 & 2.16 & 2.44 & 2.56 & 2.60 & 2.61 \\
Climatology (seasonal normal) & 1.99 & 1.99 & 1.99 & 1.99 & 1.99 & \textbf{1.99} & 1.99 \\
Damped anomaly persistence & 1.07 & 1.60 & 1.86 & 1.97 & 2.00 & 1.99 & \textbf{1.99} \\
Final model & \textbf{0.92} & \textbf{1.49} & \textbf{1.79} & \textbf{1.95} & \textbf{1.99} & 2.00 & 1.99 \\
\bottomrule
\end{tabular}

\end{table}
\fig[0.78\linewidth]{ts_13_horizon}{Forecast error versus horizon on the test period.\label{ts:fig-horizon}}
\begin{finding}
\Q At what horizon does the model stop beating simple baselines?
\F The model beats every baseline for \textbf{1 to \TSUseful{} days ahead} (MAE \TSHOne, \TSHTwo, \TSHThree, \TSHFour~°C vs
\TSHOneBase, \TSHTwoBase, \TSHThreeBase, \TSHFourBase{} for the best baseline). From day 5 all methods converge on climatology
(about \TSClim~°C): the anomaly memory has faded, exactly as the autocorrelation of Figure~\ref{ts:fig-memory} predicted.
Persistence becomes worse than the seasonal average from day \TSPersistCross{} (\TSPersistCrossMae~°C).
\end{finding}

\subsection{What drives the predicted change?}
\fig[0.78\linewidth]{ts_12_importance}{Permutation importance of each feature for the predicted one-day change (test set).\label{ts:fig-importance}}
\begin{finding}
\Q Which signals does the model use to predict tomorrow's change?
\F Today's \textbf{anomaly} (\TSImpAnom) comes first — an unusually warm day tends to cool back towards normal and vice versa.
Next is the \textbf{north–south component of the wind} (\TSImpWind): a northerly wind brings cold air, the signature of a cold
surge. Rain (\TSImpRain) and cloud cover (\TSImpCloud) follow, as they suppress daytime heating. The one-day pressure change is
used (\TSImpDp) but ranks lower than Figure~\ref{ts:fig-memory} suggests, because the wind direction carries much of the same
information.
\end{finding}

\section{Conclusions}\label{ts:conclusion}
\begin{enumerate}
\item \textbf{Machine learning beats the classical baselines, modestly}: \TSTestSkill{} lower test MAE than persistence
(\TSTestFinal{} vs \TSTestPersist~°C). Framing the target as the \emph{change} helped more than hyperparameter tuning.
\item \textbf{The hard days are cold surges.} Winter and spring are where the model gains most and where its worst misses happen;
the physical signals — northerly wind, rising pressure, today's anomaly — are what it learned to use.
\item \textbf{Skill fades after about \TSUseful{} days}, when every method converges on the seasonal average, matching the
roughly one-week memory of temperature anomalies.
\item \textbf{Data preparation lessons}: split by time, never randomly; fit the climatology on training years only; fill short
gaps by interpolating the anomaly and long gaps with climatology.
\item \textbf{Hanoi is warming} by about \TSTrend~°C per decade in this record — enough to leave a \TSResMean~°C bias in
forecasts trained on earlier years.
\end{enumerate}
\textbf{Limitations.} ERA5 is a smoothed grid-cell value rather than a thermometer in the city; real station data would be
noisier. The forecast uses the complete weather of day \emph{t}, slightly more than an evening forecast would have. Only local
variables are used, while cold surges start hundreds of kilometres to the north.

\textbf{Future work.} Add upstream features (pressure and temperature in southern China a day earlier) to see cold surges
coming; use hourly data for 24-hour forecasts; produce probabilistic forecasts (quantile regression) to warn of large drops;
compare with archived forecasts from a physics-based weather model; retrain regularly to follow the warming trend.

\begin{thebibliography}{9}
\bibitem{ts-openmeteo} Open-Meteo. \emph{Historical Weather API} (weather data by Open-Meteo.com, CC BY 4.0).
  \url{https://open-meteo.com/en/docs/historical-weather-api}
\bibitem{ts-era5} H.~Hersbach et al. \emph{The ERA5 global reanalysis}. Quarterly Journal of the Royal Meteorological Society,
  146 (730), 2020, 1999--2049. \url{https://doi.org/10.1002/qj.3803}
\bibitem{ts-sklearn} F.~Pedregosa et al. \emph{Scikit-learn: Machine Learning in Python}. Journal of Machine Learning
  Research, 12 (2011), 2825--2830.
\bibitem{ts-nb} Group G3. \emph{Notebook} \code{G3\_TimeSeries\_HanoiWeather.ipynb}. \TSNotebook
\bibitem{ts-page} Group G3. \emph{Project page (rendered notebook, report and video)}. \TSPage\ — video: \TSVideo
\end{thebibliography}


% ------------------------------------------------------------------
\sectionsflow
\chapter{Overall conclusions}\label{ch:conclusions}

\section{What the extra experiments showed}
\begin{table}[H]\centering\small
\caption{One question per extra experiment, with its numerical answer.}\label{cc:tab-questions}
\begin{tabularx}{\linewidth}{lL{5.2cm}Y}
\toprule \hd{Part} & \hd{Question} & \hd{Answer (numbers from the chapters above)} \\ \midrule
Tabular & How much do leakage and duplicates inflate the scores? & leakage: \TBLeakF{} on every metric; duplicates: F1 \TBNodupF{} instead of \TBCorrectF \\
Tabular & Does handling the class imbalance help? & recall \TBImbDefR{} → \TBImbThrR{} at threshold \TBThr; PR-AUC unchanged (\TBImbDefPR) \\
Text & Does a pretrained encoder beat TF-IDF? & test macro-F1 \TXTfidfTest{} (TF-IDF) → \TXEmbTest{} (embedding) → \TXHybTest{} (hybrid) \\
Text & Is the headline or the lead more informative? & lead \TXLeadOnly{} vs headline \TXHeadOnly; both \TXBoth \\
Time series & How should missing days be filled? & seasonal interpolation up to 7 days (\TSImpSeasSeven~°C); climatology from 14 days (\TSImpClimFourteen~°C) \\
Time series & How far ahead is the forecast useful? & beats every baseline for 1–\TSUseful{} days; converges to climatology (\TSClim~°C) after \\
\bottomrule
\end{tabularx}
\end{table}

\section{Lessons across the three parts}
\begin{enumerate}
\item \textbf{Data preparation decides whether a result can be trusted.} Leakage columns and duplicated rows in the tabular data,
ambiguous labels in the crawled news, and a random instead of chronological split in the time series would each have produced
better-looking but false numbers. In every part the honest set-up is documented and, where possible, measured.
\item \textbf{Always compare with a baseline.} A dummy classifier (accuracy \TBValDummyAcc{} but F1 0), a random guess (macro-F1
\TXRandTest) and persistence (\TSTestPersist~°C) put every score in context; the gains of the final models are real but measured
against these, not against zero.
\item \textbf{Framing matters as much as the algorithm.} Choosing the decision threshold (tabular), combining two representations
(text) and predicting the change instead of the level (time series) each helped more than hyperparameter tuning did.
\item \textbf{Errors are explainable.} Online-travel-agency bookings, overlapping news sections and winter cold surges account for
most mistakes, and each was anticipated by the exploratory analysis.
\end{enumerate}

\section{Limitations and next steps}
The tabular model was evaluated with a random split and would need a chronological test before deployment; the news dataset
covers one newspaper and one short period; the weather model uses only local variables and a reanalysis grid cell. The most
promising next steps are a time-based evaluation for the hotel data, a second crawl a few months later to measure how the text
model ages, and upstream weather features to anticipate cold surges earlier.

% ------------------------------------------------------------------
\appendix
\chapter{Links and reproduction}\label{app:links}
\begin{table}[H]\centering\small
\caption{Full addresses of every deliverable.}\label{app:tab-links}
\begin{tabularx}{\linewidth}{L{3.4cm}Y}
\toprule \hd{Item} & \hd{Address} \\ \midrule
Website (landing page) & \grpPages \\
Source repository & \grpRepo \\
Tabular — page & \TBPage \\
Tabular — notebook & \TBNotebook \\
Tabular — video & \TBVideo \\
Text — page & \TXPage \\
Text — notebook & \TXNotebook \\
Text — video & \TXVideo \\
Time series — page & \TSPage \\
Time series — notebook & \TSNotebook \\
Time series — video & \TSVideo \\
\bottomrule
\end{tabularx}
\end{table}

\textbf{Re-running the analyses.} Open a notebook from the repository in Google Colab (\emph{File → Open notebook → GitHub}) and
choose \emph{Runtime → Run all}. Each notebook loads its data from the repository and needs only the libraries pre-installed on
Colab; the text notebook additionally downloads the \code{multilingual-e5-small} model on first run.

\textbf{Rebuilding the website and reports} (from the repository root):
\begin{itemize}
\item \code{python tools/build\_pages.py} — regenerates the three project pages from the executed notebooks;
\item \code{python tools/build\_reports.py} — regenerates these PDF reports (XeLaTeX) from the notebooks' saved results.
\end{itemize}
\end{document}
